% =====================================================================
%  Bourgain-pedia digestion of:
%    J. Bourgain, T. Wolff, "A remark on gradients of harmonic functions
%    in dimension >= 3", IHES/M/89/31 (September 1989);
%    published as Colloq. Math. 60/61, No. 1, 253-260 (1990).
%    Zbl 0731.31006 -- DOI 10.4064/cm-60-61-1-253-260
%
%  Original: 11 pages (IHES preprint scan).
%  Transcription: ...-original.tex.  Ledger: ledger.md.
% =====================================================================
\documentclass[11pt,reqno]{amsart}
\usepackage[T1]{fontenc}
\usepackage{amsmath,amssymb,amsthm}
\usepackage[margin=1.05in]{geometry}
\usepackage{xcolor}
\usepackage[most]{tcolorbox}
\usepackage{enumitem}
\usepackage{hyperref}
\hypersetup{colorlinks=true,linkcolor=black,urlcolor=teal}

% ----- the four editorial environments -------------------------------
\definecolor{expcol}{HTML}{2F5D7C}
\definecolor{filcol}{HTML}{2E6B4F}
\definecolor{notcol}{HTML}{6B4E9B}
\definecolor{caucol}{HTML}{9C3B22}

\newtcolorbox{expansion}{breakable,enhanced,colback=expcol!4,colframe=expcol!55,
  boxrule=0.35pt,left=6pt,right=6pt,top=4pt,bottom=4pt,
  fonttitle=\bfseries\scriptsize,title={EXPANSION \textnormal{--- a step he compressed, written out}}}
\newtcolorbox{filled}{breakable,enhanced,colback=filcol!4,colframe=filcol!55,
  boxrule=0.35pt,left=6pt,right=6pt,top=4pt,bottom=4pt,
  fonttitle=\bfseries\scriptsize,title={FILLED \textnormal{--- an argument omitted in the original, supplied here}}}
\newtcolorbox{ournotation}{breakable,enhanced,colback=notcol!4,colframe=notcol!55,
  boxrule=0.35pt,left=6pt,right=6pt,top=4pt,bottom=4pt,
  fonttitle=\bfseries\scriptsize,title={OUR NOTATION \textnormal{--- not in the original}}}
\newtcolorbox{caution}{breakable,enhanced,colback=caucol!5,colframe=caucol!70,
  boxrule=0.8pt,left=6pt,right=6pt,top=4pt,bottom=4pt,
  fonttitle=\bfseries\scriptsize,title={CAUTION \textnormal{--- not fully justified here; OPEN-GAP in the ledger}}}

\theoremstyle{plain}
\newtheorem*{thmB}{Theorem (Bourgain--Wolff)}
\newtheorem*{lemA}{Lemma 1}
\newtheorem*{lemB}{Lemma 2}
\newtheorem*{lemC}{Lemma 3}
\newtheorem*{lemD}{Lemma 4}
\newtheorem*{corC}{Corollary}
\theoremstyle{definition}
\newtheorem*{bg}{Background}
\theoremstyle{remark}
\newtheorem*{rmk}{Remark}

\newcommand{\Rd}{\mathbb{R}^d_+}
\newcommand{\Rdm}{\mathbb{R}^{d-1}}
\newcommand{\nd}{\partial_n}
\newcommand{\he}[1]{\widehat{#1}}
\newcommand{\Ftil}{\widehat{F}_\varepsilon}
\DeclareMathOperator{\supp}{supp}
\newcommand{\Ledger}[1]{{\scriptsize\textsf{[#1]}}}

\begin{document}

\title[Gradients of harmonic functions in dimension $\ge3$ --- digestion]
      {A remark on gradients of harmonic functions in dimension $\geq 3$\\
       \normalsize a Bourgain-pedia digestion}
\author{J. Bourgain and T. Wolff}
\address{Original: IHES/M/89/31, September 1989; Colloq.\ Math.\ 60/61 (1990), 253--260}
\date{Digestion generated \today}
\maketitle

% =====================================================================
\section*{\S0. Editorial preface}
% =====================================================================

This document is the Bourgain-pedia \emph{digestion} of the Bourgain--Wolff note
\emph{A remark on gradients of harmonic functions in dimension $\ge3$}. It is that
paper at full length: every step the authors compressed is written out, every
argument they omitted is supplied, and nothing else has been changed. It is not a
survey, not a modernisation, and not a better proof. Where a modern route exists it
is confined to \S9.

Read it beside the transcription
\texttt{...-original.tex}, which reproduces the preprint as printed, including its
slips. Section and lemma numbering here is the authors'.

\medskip
\noindent\textbf{How to read the boxes.} Four kinds of material are ours, and each
is boxed so that you can see at a glance where the authors stop and we start:
\emph{expansion} (a step of theirs, written out), \emph{filled} (an argument absent
from the original), \emph{our notation}, and \emph{caution} (a step we could not
close). Unboxed prose and displays are theirs.

\medskip
\noindent\textbf{Source.} The IHES preprint (September 1989), obtained from the IHES
preprint fonds. The journal version appeared in 1990 and may differ after
refereeing; every claim here is a claim about the preprint. \Ledger{F4}

\medskip
\noindent\textbf{Five discrepancies in the original}, each verified against the scan
at 320\,dpi and each repaired below with the repair marked:
\begin{enumerate}[label=(\roman*),leftmargin=2.2em]
\item p.~4: $\int_{|x|>1}|d\rho_\varepsilon/dn|\le C\varepsilon\delta^{-p}$ is
  printed; the estimate $(*)$ yields exponent $-1$, not $-p$. Both give the
  conclusion. \Ledger{B23}
\item p.~5: the first display of Lemma 2's proof prints
  $|1+d\Ftil/dn|^p-|1+d\Ftil/dn|^p$, identically $0$; the first term must be
  $|I+d\Ftil/dn|^p$. \Ledger{C4}
\item p.~9: the final parenthesis of Lemma 4 prints $e^{+\frac12p\beta}$, which is
  $>1$ and contradicts the sentence after it; the recomputation forces
  $e^{-\frac12p\beta}$. \Ledger{D21}
\item p.~10: condition (2) prints $\varepsilon_{n+1}^{d-1}$; the use made of it
  requires $\varepsilon_{n+1}^{(d-1)/2}$. \Ledger{D23}
\item p.~10: $|Q|<K_{n+1}^{-p}e^{-\beta pn}\int_Q|du_n/dn|^p$ is printed; the
  algebra gives $e^{+\beta pn}$, and the next display on p.~11 uses $+$.
  \Ledger{D31}
\end{enumerate}
None of the five affects the theorem. Two further points needed real repair rather
than a sign change, and are treated where they arise: the lower bound for
$|1+d\he G_\varepsilon/dn|$ fails exactly at $|x|=1$ (\S2.6, \S3.3), and the final
passage from ``$f=0$ and $\partial_nf=0$'' to ``$\nabla f=0$'' is not in the paper
at all (\S7.5).

\medskip
\noindent\textbf{What we could not close.} One item, and it is the load-bearing one:
the Alexandrov--Kargaev inequality imported as reference [1] (\S2.2). Its source is
``private communication (to appear)'' and we have not obtained it. Everything else
in the paper is proved here.

% =====================================================================
\section*{\S0$'$. Notation}
% =====================================================================

\noindent\textbf{The authors' notation.}
\begin{itemize}[leftmargin=1.6em,itemsep=1pt]
\item $\Rd=\{x=(x',x_d):x'\in\Rdm,\ x_d>0\}$, boundary identified with $\Rdm$.
\item $D(0,r)\subset\Rdm$ is the open ball of radius $r$; $|E|$ is Lebesgue measure
  on $\Rdm$.
\item $Q(N)=[-\tfrac N2,\tfrac N2]^{d-1}$; in particular $Q(1)$ is the unit cube.
  $\ell(Q)$ is the side length of a cube $Q$ and $a_Q$ its centre.
\item $\he g$ is the harmonic extension of $g:\Rdm\to\mathbb R$ to $\Rd$. From \S4
  on the authors drop the hat and let $g$ denote both. \Ledger{F2}
\item $d/dn$ is the normal derivative on $\{x_d=0\}$ of a harmonic extension;
  $\nabla_T$ is differentiation in the $\Rdm$ directions.
\item $\mathcal H_n$: the $\delta_n^{-(d-1)}$ closed cubes of side $\delta_n$
  tiling $Q(1)$. $\mathcal Q_n\subset\mathcal H_n$, and
  $V_n=\bigcup\{Q:Q\in\mathcal Q_n\}$.
\item $e_d$ is the unit vector in the $x_d$ direction.
\end{itemize}

\begin{ournotation}
\begin{itemize}[leftmargin=1.6em,itemsep=1pt]
\item \textbf{$\nd$ for $d/dn$}, and we fix the convention the paper never states:
  \[\nd g(x'):=\frac{\partial \he g}{\partial x_d}(x',0^+),\]
  the derivative along the \emph{inward} normal $+e_d$. \S2.1 shows this is the
  convention the paper's estimates are written in. \Ledger{A4}
\item \textbf{$A\lesssim B$} means $A\le CB$ with $C$ depending only on $d$, on $p$,
  and on the fixed parameter $N$ once it is chosen --- never on
  $\varepsilon,\delta,\lambda,j,n$, nor on $K_n,\varepsilon_n,\delta_n$. Any other
  dependence is written out. The paper uses $\lesssim$ without ever saying this.
  \Ledger{F1}
\item \textbf{$\|g\|_p:=(\int|g|^p)^{1/p}$ for $0<p<1$.} This is \emph{not} a norm
  and obeys no triangle inequality; every step that looks like one is justified by
  Background \ref{bg:sub} instead. \Ledger{F3}
\item \textbf{$P$} for the Poisson kernel of $\Rd$ (Background \ref{bg:poisson}).
\item \textbf{$s=|x'|/\varepsilon$}, the natural variable in \S2.
\end{itemize}
\end{ournotation}

% =====================================================================
\section*{\S0$''$. Background}
% =====================================================================

Everything the paper imports, stated in full. The authors state none of these.

\begin{bg}[Poisson kernel and harmonic extension]\label{bg:poisson}
For $x'\in\Rdm$, $t>0$ put
$P(x',t)=c_d\,t\,(|x'|^2+t^2)^{-d/2}$ with $c_d=\Gamma(d/2)\pi^{-d/2}$, so that
$\int_{\Rdm}P(x',t)\,dx'=1$ for every $t>0$. For bounded continuous $g$ the
harmonic extension is $\he g(x',t)=(P(\cdot,t)*g)(x')$; it is harmonic in $\Rd$ and
$\he g(\cdot,t)\to g$ as $t\downarrow0$.
\end{bg}

\begin{bg}[Dirichlet-to-Neumann operator]\label{bg:dtn}
With $\he{g}(\xi,t)=e^{-|\xi|t}\hat g(\xi)$ on the Fourier side,
\[\nd g=\frac{\partial \he g}{\partial x_d}\Big|_{x_d=0^+}=-|\nabla'|g,
  \qquad \widehat{|\nabla'|g}(\xi)=|\xi|\hat g(\xi).\]
Since $|\nabla'|=\sum_{k=1}^{d-1}R_k\partial_k$ with $R_k$ the Riesz transforms on
$\Rdm$ (multipliers $-i\xi_k/|\xi|$),
\begin{equation}\label{eq:dtn}
  \nd g=-\sum_{k=1}^{d-1}R_k\,\partial_k g .
\end{equation}
\end{bg}

\begin{filled}
\eqref{eq:dtn} is \emph{the} reason the Riesz transforms appear in Lemma 1, and the
paper never says so: it writes ``By $L^1\to$ weak-$L^1$ and H\"older estimates for the
Riesz transforms we have\dots'' with no indication of what the Riesz transforms are
being applied to. They are applied to $\nabla_T\rho^j_\varepsilon$, and
\eqref{eq:dtn} is what converts a bound on the \emph{tangential} gradient of a
boundary function into a bound on the \emph{normal} derivative of its harmonic
extension. \Ledger{B13}
\end{filled}

\begin{bg}[Riesz transforms: the two estimates used]\label{bg:riesz}
On $\Rdm$: (i) each $R_k$ is of weak type $(1,1)$:
$|\{|R_kh|>\lambda\}|\le C_d\lambda^{-1}\|h\|_1$; (ii) for $0<\alpha<1$ each $R_k$
is bounded on the homogeneous H\"older space $C^\alpha$:
$\|R_kh\|_{C^\alpha}\le C_{d,\alpha}\|h\|_{C^\alpha}$.
\end{bg}

\begin{bg}[Representation away from the support]\label{bg:rep}
If $h\in C_c^\infty(\Rdm)$ and $x\notin\supp h$ then
$\nd h(x)=c_d'\int_{\Rdm}|x-y|^{-d}h(y)\,dy$ with $c_d'=-c_d\,(d-1)\!\int\!\cdots$;
only $|c_d'|\lesssim1$ is used. This is $\eqref{eq:dtn}$ written as a singular
integral, legitimate off the support because the kernel is then smooth.
\end{bg}

\begin{bg}[$p$-subadditivity, $0<p\le1$]\label{bg:sub}
$(a+b)^p\le a^p+b^p$ for $a,b\ge0$; consequently, for all real $a,b$,
\begin{equation}\label{eq:sub}
  \bigl||1+a|^p-|1+b|^p\bigr|\le|a-b|^p .
\end{equation}
\end{bg}

\begin{filled}
\emph{Proof of Background \ref{bg:sub}.} For $0<p\le1$ the function $t\mapsto t^p$
is concave on $[0,\infty)$ with $0^p=0$, hence subadditive: for $a,b\ge0$,
$\frac{a}{a+b}(a+b)^p+\frac{b}{a+b}(a+b)^p=(a+b)^p$ while concavity through the
origin gives $a^p\ge\frac{a}{a+b}(a+b)^p$ and $b^p\ge\frac{b}{a+b}(a+b)^p$; adding
yields $a^p+b^p\ge(a+b)^p$. For \eqref{eq:sub}, put $u=|1+a|$, $v=|1+b|$ and assume
$u\ge v$. Then $u\le v+|a-b|$ by the triangle inequality, so
$u^p\le(v+|a-b|)^p\le v^p+|a-b|^p$, i.e. $u^p-v^p\le|a-b|^p$. \hfill$\square$
\Ledger{B19}
\end{filled}

\begin{bg}[Mean value step]\label{bg:mvt}
If $|1+a|\ge c>0$ and $|1+b|\ge c>0$ then
$\bigl||1+a|^p-|1+b|^p\bigr|\le p\,c^{p-1}|a-b|$.
\end{bg}

\begin{filled}
\emph{Proof.} $t\mapsto|t|^p$ is $C^1$ away from $0$ with derivative of modulus
$p|t|^{p-1}$, which for $0<p<1$ is decreasing in $|t|$; on the segment joining
$1+a$ to $1+b$ --- which stays in $\{|t|\ge c\}$ when $1+a$ and $1+b$ have the same
sign --- the derivative is at most $pc^{p-1}$. If $1+a$ and $1+b$ have opposite signs
the segment crosses $0$; but then $|a-b|\ge2c$ and the claim follows from
$\bigl||1+a|^p-|1+b|^p\bigr|\le\max(|1+a|,|1+b|)^p$ together with
$|1+a|,|1+b|\le|a-b|+c\le\frac32|a-b|$. \hfill$\square$ \Ledger{B22}
\end{filled}

\begin{bg}[Density points]\label{bg:density}
If $E\subset\Rdm$ is measurable, a.e.\ $x\in E$ is a point of density $1$ of $E$.
If moreover $g\in C^1$ and $g\equiv0$ on $E$, then $\nabla g(x)=0$ at every density
point $x$ of $E$.
\end{bg}

\begin{filled}
\emph{Proof of the second sentence.} Let $x$ be a density point of $E$ and suppose
$\nabla g(x)=v\ne0$. Pick the unit vector $e=v/|v|$. By differentiability,
$g(x+ry)=r\langle v,y\rangle+o(r)$ uniformly for $|y|\le1$. The set
$S=\{y:|y|\le1,\ \langle v,y\rangle>\tfrac12|v|\}$ is a spherical cap of positive
measure $\kappa$, and for $r$ small $g>0$ on $x+rS$, so $(x+rS)\cap E=\emptyset$.
Then $|B(x,r)\setminus E|\ge\kappa r^{d-1}$, contradicting density $1$.
\hfill$\square$ \Ledger{D34}
\end{filled}

\begin{bg}[Maximum principle for the gradient]\label{bg:max}
If $v$ is harmonic and bounded on $\Rd$, continuous on $\overline{\Rd}$, and
$v(x)\to0$ as $|x|\to\infty$, then $\sup_{\Rd}|v|=\sup_{\Rdm}|v|$. Each partial
derivative of a harmonic function is harmonic, so the same holds for each
$\partial_jv$ provided it too is bounded and decays.
\end{bg}

% =====================================================================
\section*{\S1. The theorem}
% =====================================================================

\begin{thmB}
If $d\ge3$ there is a harmonic function $f:\Rd\to\mathbb R$ which is $C^1$ up to the
boundary and such that $f$ and $\nabla f$ vanish on a common boundary set with
positive measure.
\end{thmB}

\begin{ournotation}
The original prints ``$C'$'' for $C^1$. \Ledger{A1}
\end{ournotation}

\begin{filled}
\textbf{What the statement must mean, and one thing it must exclude.} \Ledger{A2, A3}

$f$ is harmonic in the open half space and extends to a $C^1$ function on
$\overline{\Rd}$; ``$f$ vanishes'' refers to its boundary values, and ``$\nabla f$
vanishes'' to the full gradient of that $C^1$ extension, evaluated on the boundary.
So the assertion is: there is $E\subset\Rdm$ with $|E|>0$ and
$f|_E=0$, $\nabla f|_E=0$.

As literally stated the theorem is satisfied by $f\equiv0$. The construction does
better, and nontriviality has to be read into it: we record here that the
construction produces $f\not\equiv0$, and where. In \S4 the initial datum $u_0$ is a
smooth function on $\Rdm$ vanishing on $Q(1)$; if $u_0\not\equiv0$ then, since every
later modification is supported inside $Q(1)$ (\S4.4), $f=u_0$ on
$\Rdm\setminus Q(1)$ and hence $f\not\equiv0$. \emph{So: choose $u_0$ smooth,
vanishing on $Q(1)$, and not identically zero.} The paper does not say this.
\end{filled}

\noindent\emph{Outline of proof (theirs).} Start with a function $u_0$ which
vanishes on an open subset of the boundary. Using a successive modification
(``correction theorem'') procedure, decrease the normal derivative to $0$ on a
subset with large measure. To get the theorem it is necessary to keep the original
function unchanged on a set with large measure. Thus the functions added on during
the modification procedure must have small compact support. The idea for
constructing those functions is due to A.B. Alexandrov and P. Kargaev [1].

\begin{expansion}
\textbf{The shape of the argument, in one paragraph.} The whole proof is a
contest between two budgets. Modifying $u_n$ on a cube $Q$ multiplies the local
$L^p$ size of $\partial_nu_n$ by a definite factor $e^{-2\beta}<1$ (Lemma 2) ---
that is the gain. The modification is not free: it changes $u_n$ itself on a disc
inside $Q$, and it perturbs $\partial_nu_n$ on \emph{other} cubes (Lemma 3
controls the interference). The gain is geometric in $n$; the two costs are
$\sum\varepsilon_{n+1}^{(d-1)/2}$ (measure where $f$ moved) and
$\sum K_{n+1}^{-p}$ (measure discarded from $V_n$). Condition (2) of \S7 says the
two costs together come to less than $|Q(1)|=1$, so something is left over, and
that leftover is the set $E$.
\end{expansion}

% =====================================================================
\section*{\S2. Lemma 1: the Alexandrov--Kargaev building block}
% =====================================================================

\begin{lemA}
If $p>0$ is small enough then for all sufficiently small $\varepsilon$ there is
$F_\varepsilon:\Rdm\to\mathbb R$ with
$\supp F_\varepsilon\subset D(0,\varepsilon^{1/2})$ and, $\Ftil$ denoting its
harmonic extension to $\Rd$,
\begin{equation}\label{eq:L1main}
  \int_{\Rdm}\Bigl(\bigl|1+\nd F_\varepsilon\bigr|^p-1\Bigr)dx<-\eta
\end{equation}
with $\eta>0$ independent of $\varepsilon$; moreover
\begin{equation}\label{eq:L1grad}
  |\nabla\Ftil|\lesssim\min(\varepsilon^{-d},|x|^{-d})\quad\text{on }\Rdm .
\end{equation}
\end{lemA}

\begin{ournotation}
The original prints $|x|\subset\varepsilon^{1/2}$ for $|x|<\varepsilon^{1/2}$
\Ledger{B1}, and in \eqref{eq:L1grad} prints a lower-case $\he f_\varepsilon$ where
$\Ftil$ is meant, as its own restatement on p.~4 confirms \Ledger{B2}.
\end{ournotation}

\begin{filled}
\textbf{Quantifier order.} \Ledger{B3} The statement is: there is $p_0=p_0(d)>0$
such that for every $p\in(0,p_0)$ there are $\eta=\eta(p,d)>0$ and
$\varepsilon_0=\varepsilon_0(p,d)>0$ with the following property: for every
$\varepsilon\in(0,\varepsilon_0)$ there exists $F_\varepsilon$ as described. The
point of ``$\eta$ independent of $\varepsilon$'' is that $\eta$ may depend on $p$
and $d$ but not on $\varepsilon$; \S2.7 shows where that matters. \S2.5 will force
$p<(d-1)/d$, so $p_0\le(d-1)/d$.
\end{filled}

\subsection*{2.1 The Alexandrov--Kargaev function, and the sign convention}

Alexandrov--Kargaev use the function $G_\varepsilon:\Rd\to\mathbb R$,
\begin{equation}\label{eq:G}
  G_\varepsilon(x)=\frac{-(\varepsilon+x_d)}{|x+\varepsilon e_d|^{d}} .
\end{equation}

\begin{expansion}
\textbf{$G_\varepsilon$ is harmonic on a neighbourhood of $\overline{\Rd}$.}
\Ledger{B5} By Background \ref{bg:poisson},
$P(x',t)=c_dt(|x'|^2+t^2)^{-d/2}$, so
\[G_\varepsilon(x)=-\frac1{c_d}P(x',x_d+\varepsilon),\]
i.e. $G_\varepsilon$ is $-1/c_d$ times the Poisson kernel translated \emph{down}
by $\varepsilon$. The Poisson kernel is harmonic on the open upper half space, so
$G_\varepsilon$ is harmonic on $\{x_d>-\varepsilon\}$, which contains
$\overline{\Rd}$. In particular $G_\varepsilon$ is real-analytic up to and across
the boundary $\{x_d=0\}$, and $\nd G_\varepsilon$ below is an honest classical
derivative, not a boundary limit.
\end{expansion}

\begin{expansion}
\textbf{The normal derivative, computed exactly.} Differentiating \eqref{eq:G} in
$x_d$ and setting $x_d=0$, with $r=(|x'|^2+\varepsilon^2)^{1/2}$:
\[
 \nd G_\varepsilon(x')=\frac{\partial G_\varepsilon}{\partial x_d}(x',0)
  =-r^{-d}+d\,\varepsilon^2r^{-d-2}
  =-r^{-d-2}\bigl(r^2-d\varepsilon^2\bigr).
\]
In the scaled variable $s=|x'|/\varepsilon$, so that $r=\varepsilon(s^2+1)^{1/2}$,
\begin{equation}\label{eq:dnG}
  \nd G_\varepsilon= -\,\varepsilon^{-d}\,h(s),
  \qquad h(s)=(s^2+1)^{-\frac{d+2}2}\bigl(s^2+1-d\bigr).
\end{equation}
Three features of \eqref{eq:dnG} drive everything below:
\[
 \nd G_\varepsilon(0)=(d-1)\varepsilon^{-d}>0,\qquad
 \nd G_\varepsilon=0 \text{ at } |x'|=\varepsilon\sqrt{d-1},\qquad
 \nd G_\varepsilon\sim-|x'|^{-d}\ \ (|x'|\gg\varepsilon).
\]
\end{expansion}

\begin{filled}
\textbf{Why the inward normal.} \Ledger{A4} The paper never fixes the sign of
$d/dn$, and the two choices are not interchangeable here. With the inward
convention $\nd=\partial_{x_d}|_{x_d=0}$ used above, $\nd G_\varepsilon\sim-|x'|^{-d}$
in the far field, so $1+\nd G_\varepsilon\sim1-|x'|^{-d}$ passes through $0$ near
$|x'|=1$ and is $<1$ beyond it --- and it is precisely the region
$|x'|\gtrsim1$, where $|1+\nd G_\varepsilon|^p-1\approx-p|x'|^{-d}<0$, that can
make the integral \eqref{eq:AK} negative. With the outward convention every regime
contributes a positive amount and \eqref{eq:AK} is false. So the inward normal is
the paper's convention, and we adopt it.
\end{filled}

\subsection*{2.2 The imported inequality}

As Alexandrov and Kargaev show,
\begin{equation}\label{eq:AK}
  \int_{\Rdm}\Bigl(\bigl|1+\nd G_\varepsilon\bigr|^p-1\Bigr)dx<-2M<0
\end{equation}
for $p$ below a certain critical number and $\varepsilon$ sufficiently small.

\begin{caution}
\textbf{This is the one thing in the paper we cannot verify.} \Ledger{B4}
Reference [1] is ``A.B. Aleksandrov, P. Kargaev, \emph{Private communication} (to
appear)'', and we have not obtained it; the published form is not identified
\Ledger{E4}. Everything below is therefore \emph{conditional on \eqref{eq:AK}},
which we treat as a hypothesis. Status: \textbf{OPEN-GAP}.

What can be said from here. Substituting $x=\varepsilon y$ in \eqref{eq:AK} and
using \eqref{eq:dnG}, the integral equals
$\varepsilon^{d-1}\int_{\Rdm}(|1-\varepsilon^{-d}h(|y|)|^p-1)\,dy$, and splits
into three regimes:
\begin{itemize}[leftmargin=1.4em,itemsep=1pt]
\item $|x'|\lesssim\varepsilon$: $|1+\nd G_\varepsilon|\approx(d-1)\varepsilon^{-d}$,
  contributing $O(\varepsilon^{d-1-dp})$ --- positive, $\to0$.
\item $\varepsilon\ll|x'|<1$: $|1+\nd G_\varepsilon|\approx|x'|^{-d}\gg1$, giving
  $\int_0^1(r^{-dp}-1)r^{d-2}dr=\frac{dp}{(d-1)(d-1-dp)}$ --- positive, and $O(p)$.
\item $|x'|>1$: $1+\nd G_\varepsilon\approx1-|x'|^{-d}\in(0,1)$, giving
  $\approx-p\int_{|x'|>1}|x'|^{-d}dx'$ --- \emph{negative}, and $O(p)$; together with
  the annulus around the zero at $|x'|\approx1$, where the integrand is close to
  $-1$.
\end{itemize}
So \eqref{eq:AK} asserts that at small $p$ the third regime beats the second. Both
are $O(p)$, so the inequality is a genuine computation about the constants and not
something the shape of the estimate settles. We do not reproduce it.
\end{caution}

\begin{ournotation}
The proof of Lemma 1 uses $\eta$ where \eqref{eq:AK} produced $M$; the two are
identified, $\eta:=M$. \Ledger{B6}
\end{ournotation}

\subsection*{2.3 The cut-off, and why $\supp F_\varepsilon$ is small}

Let $\delta=\varepsilon^{1/2}$ and for $j=0,1,2,\dots$ let $\psi_j:\Rdm\to\mathbb R$
be functions with
\[\supp\psi_j\subset\{2^{j-1}\delta\le|x|\le2^{j+1}\delta\},\qquad
  \sum_j\psi_j=1\ \text{ for }|x|>\delta,\qquad
  |\nabla^k\psi_j|\lesssim(2^j\delta)^{-k}.\]
Put $\psi=\sum_j\psi_j$, $\rho^j_\varepsilon=\psi_jG_\varepsilon$,
$\rho_\varepsilon=\sum_j\rho^j_\varepsilon=\psi G_\varepsilon$, and
\[F_\varepsilon=G_\varepsilon-\rho_\varepsilon=(1-\psi)G_\varepsilon .\]

\begin{filled}
\textbf{Such a family exists.} \Ledger{B7} Fix $\chi\in C^\infty_c(\mathbb R)$ with
$\supp\chi\subset[\frac12,2]$, $\chi\ge0$, and $\sum_{j\in\mathbb Z}\chi(2^{-j}t)=1$
for $t>0$ (a standard dyadic partition: take any $\theta\in C^\infty_c$ with
$\theta=1$ on $[0,1]$, $\theta=0$ on $[2,\infty)$, and set
$\chi(t)=\theta(t)-\theta(2t)$). Put $\psi_j(x)=\chi\!\left(|x|/(2^j\delta)\right)$
for $j\ge0$. Then $\supp\psi_j\subset\{2^{j-1}\delta\le|x|\le2^{j+1}\delta\}$;
$\sum_{j\ge0}\psi_j(x)=1$ once $|x|>\delta$, because the omitted terms $j<0$ are
supported in $\{|x|<\delta\}$; and $|\nabla^k\psi_j|\le\|\chi^{(k)}\|_\infty
(2^j\delta)^{-k}$ by the chain rule.
\end{filled}

\begin{expansion}
\textbf{$\supp F_\varepsilon\subset D(0,\varepsilon^{1/2})$, and $F_\varepsilon$ is
smooth.} \Ledger{B8} With the $\psi_j$ above, $\psi=0$ on $\{|x|<\delta/2\}$ (every
$\psi_j$, $j\ge0$, vanishes there) and $\psi=1$ on $\{|x|>\delta\}$. Hence
$1-\psi$ is supported in $\overline{D(0,\delta)}$, and since
$F_\varepsilon=(1-\psi)G_\varepsilon$,
\[\supp F_\varepsilon\subset\overline{D(0,\delta)}=\overline{D(0,\varepsilon^{1/2})},\]
which is the claim in Lemma 1. Smoothness: on the boundary
$G_\varepsilon(x',0)=-\varepsilon(|x'|^2+\varepsilon^2)^{-d/2}$ is real-analytic
(the denominator never vanishes, as $\varepsilon>0$), and $1-\psi$ is smooth; so
$F_\varepsilon\in C_c^\infty(\Rdm)$. Note in passing that $F_\varepsilon$ agrees
with $G_\varepsilon$ on $\{|x|<\delta/2\}$: the cut-off removes the far field, not
the singular core.
\end{expansion}

\subsection*{2.4 Estimates on $\rho_\varepsilon$}

Let $\nabla_T$ be differentiation in the $\Rdm$ directions. Then
$|\nabla_T^kG_\varepsilon|\lesssim\varepsilon|x|^{-(d+k)}$ for $k=0,1,2$; therefore
$|\nabla_T^k\rho^j_\varepsilon|\lesssim\varepsilon(2^j\delta)^{-(d+k)}$, whence
\[\|\nabla_T\rho^j_\varepsilon\|_1\lesssim\varepsilon(2^j\delta)^{-2},
  \qquad
  \|\nabla_T\rho^j_\varepsilon\|_{C^\alpha}\lesssim\varepsilon(2^j\delta)^{-(d+1+\alpha)} .\]

\begin{expansion}
\textbf{The three estimates, done.} \Ledger{B9, B10, B11, B12}
On the boundary, $G_\varepsilon(x',0)=-\varepsilon(|x'|^2+\varepsilon^2)^{-d/2}$.
\begin{itemize}[leftmargin=1.4em,itemsep=2pt]
\item $k=0$: $|G_\varepsilon|=\varepsilon(|x'|^2+\varepsilon^2)^{-d/2}
  \le\varepsilon|x'|^{-d}$, with constant $1$.
\item $k=1$: $\nabla_TG_\varepsilon
  =d\,\varepsilon\,x'(|x'|^2+\varepsilon^2)^{-\frac d2-1}$, so
  $|\nabla_TG_\varepsilon|\le d\,\varepsilon|x'|\cdot|x'|^{-d-2}
  =d\,\varepsilon|x'|^{-(d+1)}$.
\item $k=2$: differentiating once more produces two terms, of sizes
  $d\varepsilon(|x'|^2+\varepsilon^2)^{-\frac d2-1}$ and
  $d(d+2)\varepsilon|x'|^2(|x'|^2+\varepsilon^2)^{-\frac d2-2}$, both
  $\le C_d\,\varepsilon|x'|^{-(d+2)}$.
\end{itemize}
Note the paper writes ``$\le$'' where the constant depends on $d$; we write
$\lesssim$.

For $\rho^j_\varepsilon=\psi_jG_\varepsilon$, Leibniz gives
$\nabla_T^k\rho^j_\varepsilon=\sum_{i\le k}\binom ki
 (\nabla_T^i\psi_j)(\nabla_T^{k-i}G_\varepsilon)$, and on $\supp\psi_j$ we have
$|x|\sim2^j\delta$, so
\[|\nabla_T^k\rho^j_\varepsilon|
 \lesssim\sum_{i\le k}(2^j\delta)^{-i}\cdot\varepsilon(2^j\delta)^{-(d+k-i)}
 \lesssim\varepsilon(2^j\delta)^{-(d+k)} .\]
For the $L^1$ bound, $\supp\nabla_T\rho^j_\varepsilon$ lies in an annulus of measure
$\lesssim(2^j\delta)^{d-1}$, so
\[\|\nabla_T\rho^j_\varepsilon\|_1
 \lesssim\varepsilon(2^j\delta)^{-(d+1)}\cdot(2^j\delta)^{d-1}
 =\varepsilon(2^j\delta)^{-2}. \]
For the H\"older bound, fix once and for all an $\alpha\in(0,1)$ (any choice works;
$\alpha=\frac12$ will do). A function with
$|\nabla_T\rho^j|\lesssim\varepsilon(2^j\delta)^{-(d+1)}$ and
$|\nabla_T^2\rho^j|\lesssim\varepsilon(2^j\delta)^{-(d+2)}$ satisfies
$\|\nabla_T\rho^j\|_{C^\alpha}\lesssim
 \|\nabla_T\rho^j\|_\infty^{1-\alpha}\|\nabla^2_T\rho^j\|_\infty^{\alpha}
 \lesssim\varepsilon(2^j\delta)^{-(d+1+\alpha)}$, which is the stated bound.
\end{expansion}

By \eqref{eq:dtn} and Background \ref{bg:riesz} applied to
$h=\partial_k\rho^j_\varepsilon$,
\begin{equation}\label{eq:weak}
  \Bigl|\Bigl\{x:\bigl|\nd\rho^j_\varepsilon(x)\bigr|>\lambda\Bigr\}\Bigr|
   \lesssim\varepsilon(2^j\delta)^{-2}\lambda^{-1},
\end{equation}
\begin{equation}\label{eq:hold}
  \bigl\|\nd\rho^j_\varepsilon\bigr\|_{C^\alpha}
   \lesssim\varepsilon(2^j\delta)^{-(d+1+\alpha)} .
\end{equation}

\begin{expansion}
\textbf{Near field: $|\nd\rho^j_\varepsilon(x)|\lesssim\varepsilon(2^j\delta)^{-(d+1)}$
for $|x|<2\cdot2^{j+1}\delta$.} \Ledger{B15} Write $R=2\cdot2^{j+1}\delta$,
$T=\varepsilon(2^j\delta)^{-(d+1)}$, and $g=\nd\rho^j_\varepsilon$.

By \eqref{eq:weak} with $\lambda=T$, the set $\{|g|>T\}$ has measure
$\lesssim\varepsilon(2^j\delta)^{-2}T^{-1}=(2^j\delta)^{d-1}$. The ball $D(0,R)$ has
measure $c_{d}R^{d-1}=c_d\,4^{d-1}(2^j\delta)^{d-1}$. So if the implicit constant in
\eqref{eq:weak} is $C_1$, then applying \eqref{eq:weak} at level $\lambda=KT$
instead, with $K$ a constant chosen so that $C_1K^{-1}<c_d4^{d-1}$, gives
\[|\{|g|>KT\}\cap D(0,R)|<|D(0,R)|,\]
so there exists $x_0\in D(0,R)$ with $|g(x_0)|\le KT$. That is the ``\emph{some} $x$''
of the original.

To pass to \emph{all} $x\in D(0,R)$, use \eqref{eq:hold}: for any $x\in D(0,R)$,
\[|g(x)|\le|g(x_0)|+\|g\|_{C^\alpha}|x-x_0|^\alpha
 \le KT+C_2\,\varepsilon(2^j\delta)^{-(d+1+\alpha)}(2R)^\alpha .\]
Since $2R=8\cdot2^j\delta$, the second term is
$C_2 8^\alpha\varepsilon(2^j\delta)^{-(d+1)}=C_2 8^\alpha T$. Hence
$|g(x)|\le(K+C_28^\alpha)T\lesssim T$ throughout $D(0,R)$, as claimed. The
constants depend only on $d$ and $\alpha$, both fixed.
\end{expansion}

Furthermore, if $|x|>2\cdot2^{j+1}\delta$ then, by Background \ref{bg:rep},
\begin{equation}\label{eq:far}
 \bigl|\nd\rho^j_\varepsilon(x)\bigr|
  =\Bigl|c_d'\!\!\int_{|y|<2^{j+1}\delta}\!\!|x-y|^{-d}\rho^j_\varepsilon(y)\,dy\Bigr|
  \lesssim\varepsilon(2^j\delta)^{-1}|x|^{-d},
\end{equation}
using $|\rho^j_\varepsilon|\lesssim\varepsilon(2^j\delta)^{-d}$.

\begin{expansion}
\textbf{Where \eqref{eq:far} comes from.} \Ledger{B16} Since
$|x|>2\cdot2^{j+1}\delta$ and $|y|<2^{j+1}\delta$ we have $|x-y|\ge\frac12|x|$, so
$|x-y|^{-d}\le2^d|x|^{-d}$ and the integral is at most
$2^d|x|^{-d}\|\rho^j_\varepsilon\|_1$. Now
$\|\rho^j_\varepsilon\|_1\lesssim\varepsilon(2^j\delta)^{-d}\cdot(2^j\delta)^{d-1}
 =\varepsilon(2^j\delta)^{-1}$, the annulus again having measure
$\lesssim(2^j\delta)^{d-1}$. This gives \eqref{eq:far}. The representation itself is
legitimate because $x$ lies outside $\supp\rho^j_\varepsilon$, where the kernel of
$|\nabla'|$ is smooth.
\end{expansion}

In other words
\begin{equation}\label{eq:combined}
 \bigl|\nd\rho^j_\varepsilon\bigr|
  \lesssim\varepsilon(2^j\delta)^{-1}\min\!\bigl(|x|^{-d},(2^j\delta)^{-d}\bigr)
  \quad\text{on }\Rdm,
\end{equation}
and summing over $j$,
\begin{equation}\label{eq:star}\tag{$*$}
 \bigl|\nd\rho_\varepsilon\bigr|
  \lesssim\varepsilon\,\delta^{-1}\min\!\bigl(\delta^{-d},|x|^{-d}\bigr).
\end{equation}

\begin{expansion}
\textbf{The summation.} \Ledger{B17} Fix $x$ and let $j_0$ be the least $j$ with
$2^j\delta\ge|x|$. For $j<j_0$ the minimum in \eqref{eq:combined} is $|x|^{-d}$ and
the terms are $\varepsilon(2^j\delta)^{-1}|x|^{-d}$, a geometric series in $2^{-j}$
summing to $\lesssim\varepsilon\delta^{-1}|x|^{-d}$ (dominated by $j=0$). For
$j\ge j_0$ the minimum is $(2^j\delta)^{-d}$ and the terms are
$\varepsilon(2^j\delta)^{-(d+1)}$, again geometric, summing to
$\lesssim\varepsilon(2^{j_0}\delta)^{-(d+1)}\le\varepsilon\delta^{-1}|x|^{-d}$
whenever $|x|\ge\delta$. For $|x|<\delta$ only $j\ge j_0=0$ occurs and the sum is
$\lesssim\varepsilon\delta^{-(d+1)}=\varepsilon\delta^{-1}\delta^{-d}$. Both cases
are \eqref{eq:star}.
\end{expansion}

\subsection*{2.5 The main integral}

Then for small $p$ and appropriate $\eta$, using
$\nd F_\varepsilon=\nd G_\varepsilon-\nd\rho_\varepsilon$,
\begin{align}
 \int\Bigl(\bigl|1+\nd F_\varepsilon\bigr|^p-1\Bigr)
 &=\underbrace{\int\Bigl(\bigl|1+\nd G_\varepsilon\bigr|^p-1\Bigr)}_{<\,-2\eta
   \ \text{by \eqref{eq:AK}}}
  +\int\Bigl(\bigl|1+\nd F_\varepsilon\bigr|^p-\bigl|1+\nd G_\varepsilon\bigr|^p\Bigr)
  \notag\\
 &\le-2\eta
  +\underbrace{\int_{|x|<1}\Bigl|\,\bigl|1+\nd F_\varepsilon\bigr|^p
    -\bigl|1+\nd G_\varepsilon\bigr|^p\Bigr|}_{=:\,\mathrm{I}}\notag\\
 &\qquad\qquad
  +\underbrace{\int_{|x|>1}\Bigl(\bigl|1+\nd F_\varepsilon\bigr|^p
    -\bigl|1+\nd G_\varepsilon\bigr|^p\Bigr)}_{=:\,\mathrm{II}} .
  \label{eq:split}
\end{align}

\noindent\textbf{The term I.} By Background \ref{bg:sub},
\[\mathrm I\le\int_{|x|\le1}\bigl|\nd\rho_\varepsilon\bigr|^p ,\]
and $(*)$ implies this goes to $0$ with $\varepsilon$.

\begin{expansion}
\textbf{It goes to $0$ --- and this is where $p$ is constrained.} \Ledger{B20}
By \eqref{eq:star}, $|\nd\rho_\varepsilon|\lesssim\varepsilon\delta^{-1}
\min(\delta^{-d},|x|^{-d})$, and $\varepsilon\delta^{-1}=\varepsilon^{1/2}$ since
$\delta=\varepsilon^{1/2}$. Split at $|x|=\delta$.
\begin{itemize}[leftmargin=1.4em,itemsep=2pt]
\item $|x|<\delta$: the integrand is
  $\lesssim(\varepsilon^{1/2}\delta^{-d})^p$ over a set of measure
  $\lesssim\delta^{d-1}$, giving
  $\varepsilon^{p/2}\delta^{-dp+d-1}
   =\varepsilon^{\frac p2}\varepsilon^{\frac{-dp+d-1}{2}}
   =\varepsilon^{\frac{(d-1)(1-p)}{2}}\longrightarrow0$ for $p<1$.
\item $\delta<|x|\le1$: the integrand is
  $\lesssim\varepsilon^{p/2}|x|^{-dp}$, and
  \[\int_\delta^1 r^{-dp}\,r^{d-2}\,dr
    =\int_\delta^1 r^{\,d-2-dp}\,dr,\]
  which stays bounded as $\delta\downarrow0$ \emph{iff} $d-2-dp>-1$, i.e.
  \[\boxed{\,p<\tfrac{d-1}{d}\,}\]
  and then the whole piece is $\lesssim\varepsilon^{p/2}\to0$.
\end{itemize}
The paper says only ``for small $p$''. The explicit constraint is
$p<\frac{d-1}d$, which for $d\ge3$ means $p<\frac23$; combined with $p<1$ from
Background \ref{bg:sub} the binding constraint is $p<\frac{d-1}{d}$. Together with
whatever ``below a certain critical number'' means in \eqref{eq:AK}, this fixes
$p_0(d)$.
\end{expansion}

\noindent\textbf{The term II.} When $|x|>1$ we have a lower bound for
$|1+\nd G_\varepsilon|$, and $|\nd\rho_\varepsilon|$ is small; it follows by the
mean value theorem that
$\bigl||1+\nd F_\varepsilon|^p-|1+\nd G_\varepsilon|^p\bigr|
 \le C|\nd\rho_\varepsilon|$, and therefore
\[\mathrm{II}\le C\int_{|x|>1}\bigl|\nd\rho_\varepsilon\bigr|
  \le C\varepsilon\delta^{-p},\]
which again goes to $0$ with $\varepsilon$. That proves \eqref{eq:L1main}.

\subsection*{2.6 Two repairs in the term II}

\begin{caution}[downgraded: repaired below]
\textbf{The lower bound fails exactly at $|x|=1$.} \Ledger{B21} By
\eqref{eq:dnG}, in the far field $\nd G_\varepsilon\approx-|x|^{-d}$, so
\[1+\nd G_\varepsilon\approx1-|x|^{-d},\]
which \emph{vanishes at $|x|\approx1$} and is only $\approx d(|x|-1)$ just beyond
it. So ``we have a lower bound for $|1+\nd G_\varepsilon|$ when $|x|>1$'' is not
true as stated: the mean value theorem is being applied across the very point where
the function it is applied to has a zero. The following repair costs nothing.
\end{caution}

\begin{filled}
\textbf{Repair: split $\{|x|>1\}$ at $|x|=2$.}
\emph{(a) On $\{1<|x|<2\}$}, do not use the mean value theorem. Background
\ref{bg:sub} gives directly
\[\int_{1<|x|<2}\Bigl|\,\bigl|1+\nd F_\varepsilon\bigr|^p
   -\bigl|1+\nd G_\varepsilon\bigr|^p\Bigr|
  \le\int_{1<|x|<2}\bigl|\nd\rho_\varepsilon\bigr|^p
  \lesssim\bigl(\varepsilon^{1/2}\bigr)^{p}\longrightarrow0,\]
using $|\nd\rho_\varepsilon|\lesssim\varepsilon^{1/2}|x|^{-d}\le\varepsilon^{1/2}$
there and the finite measure of the annulus. This is the same computation as for
term I, on a region of measure $O(1)$.

\emph{(b) On $\{|x|>2\}$} the lower bound is genuine: $|x|^{-d}\le2^{-d}\le\frac12$,
so $1+\nd G_\varepsilon\ge1-2^{-d}\ge\frac12$; and
$|\nd\rho_\varepsilon|\lesssim\varepsilon^{1/2}|x|^{-d}\le\varepsilon^{1/2}2^{-d}$,
so also $1+\nd F_\varepsilon\ge\frac12-\varepsilon^{1/2}\ge\frac14$ for
$\varepsilon$ small. Background \ref{bg:mvt} with $c=\frac14$ then gives
$\bigl||1+\nd F_\varepsilon|^p-|1+\nd G_\varepsilon|^p\bigr|
 \le p4^{1-p}\,|\nd\rho_\varepsilon|$, so
\[\int_{|x|>2}\bigl|\cdots\bigr|\le C\int_{|x|>2}\bigl|\nd\rho_\varepsilon\bigr| .\]
The repaired split changes nothing downstream: both pieces $\to0$ with
$\varepsilon$, which is all \eqref{eq:split} needs.
\end{filled}

\begin{expansion}
\textbf{The exponent in $C\varepsilon\delta^{-p}$.} \Ledger{B23} What $(*)$ yields
is
\[\int_{|x|>2}\bigl|\nd\rho_\varepsilon\bigr|
 \lesssim\varepsilon\delta^{-1}\int_{|x|>2}|x|^{-d}dx
 =c\,\varepsilon\delta^{-1}=c\,\varepsilon^{1/2},\]
since $\int_{|x|>2}|x|^{-d}dx=c\int_2^\infty r^{-d}r^{d-2}dr=c\int_2^\infty
r^{-2}dr<\infty$ in $\Rdm$. The original prints the exponent $-p$ rather than $-1$;
as $\delta<1$ and $p<1$, $\delta^{-p}<\delta^{-1}$, so the printed bound is the
\emph{stronger} of the two and is not what $(*)$ delivers. Since both
$\varepsilon\delta^{-1}=\varepsilon^{1/2}$ and
$\varepsilon\delta^{-p}=\varepsilon^{1-p/2}$ tend to $0$, the conclusion is
unaffected. We use $\varepsilon^{1/2}$.
\end{expansion}

\subsection*{2.7 The gradient bound \eqref{eq:L1grad}}

To prove $|\nabla\Ftil|\lesssim\min(\varepsilon^{-d},|x|^{-d})$: this estimate is
obvious for $|\nabla G_\varepsilon|$ and for $|\nabla_T\rho_\varepsilon|$; for
$|\nd\rho_\varepsilon|$ it follows from $(*)$.

\begin{expansion}
\textbf{``Obvious'', done.} \Ledger{B24} Since
$\nabla\Ftil=\nabla G_\varepsilon-\nabla\rho_\varepsilon$ and
$\nabla=(\nabla_T,\nd)$, there are three pieces.

\emph{(i) $\nabla G_\varepsilon$.} The normal component is \eqref{eq:dnG}:
$|\nd G_\varepsilon|=\varepsilon^{-d}|h(s)|$ with $h$ bounded, so
$|\nd G_\varepsilon|\le(d-1)\varepsilon^{-d}$ always, while for $|x|\gg\varepsilon$,
$|h(s)|\approx s^{-d}$ gives $|\nd G_\varepsilon|\approx|x|^{-d}$. Hence
$|\nd G_\varepsilon|\lesssim\min(\varepsilon^{-d},|x|^{-d})$ --- this is exactly the
shape asserted, and it is the term that forces the $\varepsilon^{-d}$ in
\eqref{eq:L1grad}. The tangential component is
$|\nabla_TG_\varepsilon|\lesssim\varepsilon|x|^{-(d+1)}$ from \S2.4, which is
smaller.

\emph{(ii) $\nabla_T\rho_\varepsilon$.} On $\supp\psi_j$, \S2.4 gives
$|\nabla_T\rho^j_\varepsilon|\lesssim\varepsilon(2^j\delta)^{-(d+1)}
 \approx\varepsilon|x|^{-(d+1)}$; the $\psi_j$ have bounded overlap, so
$|\nabla_T\rho_\varepsilon|\lesssim\varepsilon|x|^{-(d+1)}$, and since
$\rho_\varepsilon$ is supported in $\{|x|>\delta/2\}$ this is
$\le2\varepsilon\delta^{-1}|x|^{-d}=2\varepsilon^{1/2}|x|^{-d}\le|x|^{-d}$.

\emph{(iii) $\nd\rho_\varepsilon$.} By $(*)$,
$|\nd\rho_\varepsilon|\lesssim\varepsilon^{1/2}\min(\delta^{-d},|x|^{-d})$. The
second entry is already $\le|x|^{-d}$; for the first,
$\varepsilon^{1/2}\delta^{-d}=\varepsilon^{\frac{1-d}2}\le\varepsilon^{-d}$
because $\frac{d-1}2\le d$ for $d\ge1$.

Adding, $|\nabla\Ftil|\lesssim\min(\varepsilon^{-d},|x|^{-d})$. \hfill$\square$
\end{expansion}

% =====================================================================
\section*{\S3. Lemma 2: the correction step on a cube}
% =====================================================================

\begin{lemB}
If $N$ is large enough there is a constant $\beta=\beta(N)>0$ such that: if
$Q\subset\Rdm$ is a cube, $a_Q$ its centre and $I:Q\to\mathbb R$ a function such
that
\begin{equation}\label{eq:L2hyp}
  N^{d-1}\,|I(a_Q)|^{-1}\sup_{x\in Q}|I(x)-I(a_Q)|
\end{equation}
is sufficiently small (independently of $N$ and $\varepsilon$), then
\begin{equation}\label{eq:L2}
 \Bigl(\int_Q\Bigl|I(x)+\nd F_\varepsilon\bigl(N\ell(Q)^{-1}(x-a_Q)\bigr)\Bigr|^p dx
 \Bigr)^{1/p}<e^{-2\beta}\,|I(a_Q)|\,|Q|^{1/p}.
\end{equation}
\end{lemB}

\begin{ournotation}
The original prints $\beta\subset\beta(N)$ for $\beta=\beta(N)$. \Ledger{C1}
\end{ournotation}

\begin{filled}
\textbf{Order of quantifiers.} \Ledger{C2} Read: there are $N_0(d,p)$ and an
absolute threshold $\tau=\tau(d,p)>0$ such that for every $N\ge N_0$ there is
$\beta=\beta(N,d,p)>0$ such that for all $\varepsilon<\varepsilon_0(p,d)$ (as in
Lemma 1), all cubes $Q$, and all $I$ with \eqref{eq:L2hyp} $<\tau$, the conclusion
\eqref{eq:L2} holds. The parenthetical ``independently of $N$ and $\varepsilon$''
refers to $\tau$, not to $\beta$: $\beta$ certainly depends on $N$.
\end{filled}

\subsection*{3.1 Reduction to $Q=Q(N)$ and $I(a_Q)=1$}

We can assume $Q=Q(N)$. It is then equivalent to show that if $N^{d-1}\|I-1\|_\infty$
is sufficiently small then
\begin{equation}\label{eq:L2red}
 \int_{Q(N)}\Bigl(\bigl|I+\nd F_\varepsilon\bigr|^p-1\Bigr)\le-\eta
\end{equation}
for suitable $\eta>0$.

\begin{expansion}
\textbf{The two reductions.} \Ledger{C3}
\emph{Scaling.} The map $\Phi(x)=N\ell(Q)^{-1}(x-a_Q)$ carries $Q$ onto $Q(N)$ and
$a_Q$ to $0$; it has constant Jacobian $(N/\ell(Q))^{d-1}$, and
$|Q|=\ell(Q)^{d-1}$, $|Q(N)|=N^{d-1}$, so $|Q(N)|/|Q|=(N/\ell(Q))^{d-1}$. Writing
$\tilde I=I\circ\Phi^{-1}$ and substituting $y=\Phi(x)$,
\[\int_Q\bigl|I(x)+\nd F_\varepsilon(\Phi(x))\bigr|^pdx
 =\frac{|Q|}{|Q(N)|}\int_{Q(N)}\bigl|\tilde I(y)+\nd F_\varepsilon(y)\bigr|^pdy .\]
Since $|Q|^{1/p}$ appears on the right of \eqref{eq:L2}, the inequality is
equivalent to
$\frac1{|Q(N)|}\int_{Q(N)}|\tilde I+\nd F_\varepsilon|^p<e^{-2\beta p}|\tilde I(0)|^p$,
which involves $Q$ only through $\tilde I$; and \eqref{eq:L2hyp} is invariant under
$\Phi$ because it is built from the values of $I$ alone.

\emph{Normalisation.} Both sides of the resulting inequality are homogeneous of
degree $p$ in $(\tilde I,\nd F_\varepsilon)$ jointly --- but $\nd F_\varepsilon$ is
fixed, so we cannot rescale it. What we can do is divide by $|\tilde I(0)|$: it is
enough to prove the statement when $\tilde I(0)=1$, replacing $F_\varepsilon$ by
$F_\varepsilon/|\tilde I(0)|$ --- and here the argument uses that Lemma 1 is applied
to a \emph{fixed} $F_\varepsilon$ while $I$ varies. In the application (\S6) $I$ is
$\partial_nu_n(a_Q)$ plus a small perturbation, and the factor $\partial_nu_n(a_Q)$
is exactly the coefficient multiplying $F_{\varepsilon_{n+1}}$ in the definition of
$u_{n+1}$, so the two scalings match. With $\tilde I(0)=1$, \eqref{eq:L2hyp} becomes
$N^{d-1}\|\tilde I-1\|_{\infty}$, and \eqref{eq:L2} becomes \eqref{eq:L2red} with
$e^{-2\beta p}=1-\eta/N^{d-1}$; see \S3.4.
\end{expansion}

\subsection*{3.2 The estimate}

This last inequality is true since
\begin{align}
 \int_{Q(N)}\Bigl(\bigl|I+\nd F_\varepsilon\bigr|^p
   -\bigl|1+\nd F_\varepsilon\bigr|^p\Bigr)dx
 &\le\int_{Q(N)\setminus D(0,1)}+\int_{D(0,1)} \notag\\
 &\le C\!\int_{Q(N)\setminus D(0,1)}\!\!|I-1|\,dx
   +\int_{D(0,1)}|I-1|^p dx \label{eq:L2chain}
\end{align}
(because $1+\nd F_\varepsilon$ is bounded away from $0$ when $|x|>1$)
\[\le C\|I-1\|_\infty N^{d-1}+C\|I-1\|_\infty^{\,p},\]
which is small. It remains to apply Lemma 1.

\begin{expansion}
\textbf{Three corrections in that chain.}
\begin{itemize}[leftmargin=1.4em,itemsep=3pt]
\item \Ledger{C4} The original prints both terms of the first line as
  $|1+\nd F_\varepsilon|^p$, so the left side is identically $0$. The first must be
  $|I+\nd F_\varepsilon|^p$: the whole point is to compare the $I$-version with the
  $1$-version, which is what Lemma 1 controls. We have restored $I$.
\item \Ledger{C6} The original's last line prints $C\|I-1\|_\infty$ for the second
  term. That term is $\int_{D(0,1)}|I-1|^p\le|D(0,1)|\cdot\|I-1\|_\infty^p$, so the
  exponent is $p$, not $1$. Since $\|I-1\|_\infty<1$ and $p<1$ we have
  $\|I-1\|_\infty^p>\|I-1\|_\infty$, so the printed bound is the smaller one and is
  not what the line before yields. Harmless: the hypothesis
  $N^{d-1}\|I-1\|_\infty<\tau$ with $N\ge1$ forces $\|I-1\|_\infty<\tau$, hence
  $\|I-1\|_\infty^p<\tau^p$, and $\tau$ is at our disposal.
\item The first term: $\int_{Q(N)\setminus D(0,1)}|I-1|\le\|I-1\|_\infty|Q(N)|
  =\|I-1\|_\infty N^{d-1}$, which is where the $N^{d-1}$ in \eqref{eq:L2hyp} comes
  from --- and it is the only place. This is why the hypothesis carries the factor
  $N^{d-1}$ and not, say, $N^{d}$.
\end{itemize}
\end{expansion}

\subsection*{3.3 The same lower-bound repair}

\begin{caution}[downgraded: repaired below]
The parenthetical ``$1+\nd F_\varepsilon$ is bounded away from $0$ when $|x|>1$''
has exactly the defect of \S2.6: by \eqref{eq:dnG} and \S2.4,
$\nd F_\varepsilon=\nd G_\varepsilon-\nd\rho_\varepsilon\approx-|x|^{-d}$ in the far
field, so $1+\nd F_\varepsilon$ has a zero near $|x|=1$. \Ledger{C5}
\end{caution}

\begin{filled}
\textbf{Repair.} Split $Q(N)\setminus D(0,1)$ into the annulus
$A=\{1<|x|<2\}$ and $Q(N)\setminus D(0,2)$.
On $A$, use Background \ref{bg:sub}: the integrand is $\le|I-1|^p$, so this piece
contributes $\le|A|\cdot\|I-1\|_\infty^p\le C\|I-1\|_\infty^p$ --- the same shape as
the $D(0,1)$ term, and absorbed into it.
On $Q(N)\setminus D(0,2)$ the bound $|1+\nd F_\varepsilon|\ge\frac14$ of \S2.6(b)
holds, so Background \ref{bg:mvt} applies and gives $C|I-1|$ pointwise, hence
$C\|I-1\|_\infty N^{d-1}$.
The final bound is therefore
$C\|I-1\|_\infty N^{d-1}+C\|I-1\|_\infty^{\,p}$, exactly as in \S3.2, and the
argument proceeds unchanged.
\end{filled}

\subsection*{3.4 Applying Lemma 1, and what $\beta$ is}

\begin{filled}
\textbf{``It remains to apply lemma 1'' --- the tail.} \Ledger{C7} Lemma 1 controls
$\int_{\Rdm}(|1+\nd F_\varepsilon|^p-1)<-\eta_1$, an integral over the whole
hyperplane, whereas \eqref{eq:L2red} is over $Q(N)$. The difference is
\[\int_{\Rdm\setminus Q(N)}\Bigl(\bigl|1+\nd F_\varepsilon\bigr|^p-1\Bigr),\]
and $\Rdm\setminus Q(N)\subset\{|x|>N/2\}$. There
$|\nd F_\varepsilon|\lesssim|x|^{-d}\le(N/2)^{-d}\le\frac12$ for $N$ large, so by
Background \ref{bg:mvt} with $c=\frac12$,
$\bigl||1+\nd F_\varepsilon|^p-1\bigr|\le p2^{1-p}|\nd F_\varepsilon|
 \lesssim|x|^{-d}$, and
\[\Bigl|\int_{\Rdm\setminus Q(N)}\Bigr|\lesssim\int_{|x|>N/2}|x|^{-d}dx
 =c\,N^{-1}.\]
Choosing $N$ large enough that $cN^{-1}<\eta_1/2$ gives
\[\int_{Q(N)}\Bigl(\bigl|1+\nd F_\varepsilon\bigr|^p-1\Bigr)<-\tfrac{\eta_1}2 .\]
Adding \eqref{eq:L2chain}, whose right-hand side is $<\eta_1/4$ once
$N^{d-1}\|I-1\|_\infty<\tau$ with $\tau$ small, yields \eqref{eq:L2red} with
$\eta=\eta_1/4$. \emph{Note the order}: $\eta_1$ comes from Lemma 1 and depends on
$p,d$ only; $N$ is then chosen; then $\tau$; and only then $\beta$.
\end{filled}

\begin{expansion}
\textbf{Where $e^{-2\beta}$ comes from, and what it depends on.} \Ledger{C8}
\eqref{eq:L2red} says $\int_{Q(N)}|I+\nd F_\varepsilon|^p\le|Q(N)|-\eta$. Hence
\[\Bigl(\int_{Q(N)}\bigl|I+\nd F_\varepsilon\bigr|^p\Bigr)^{1/p}
 \le\bigl(1-\eta/N^{d-1}\bigr)^{1/p}|Q(N)|^{1/p},\]
so \eqref{eq:L2} holds with
\[\boxed{\ \beta=\beta(N)=-\frac1{2p}\log\bigl(1-\eta N^{-(d-1)}\bigr)>0\ }\]
which depends on $p$, $d$ and $N$ --- and on nothing else. This is the $\beta$ that
governs the geometric decay of the whole construction, and it is fixed once and
for all as soon as $N$ is. \hfill$\square$
\end{expansion}

% =====================================================================
\section*{\S4. The recursive construction}
% =====================================================================

For a sufficiently large $N$ to be determined, choose $\beta$ according to Lemma 2.
Also fix sequences $K_n\to\infty$, $\varepsilon_n\searrow0$ (to be determined).
Constants denoted $C$ will be independent of these choices. Let $u_0$ be a smooth
function on $\Rdm$ which vanishes on $Q(1)$.

\begin{filled}
\textbf{$u_0$ vanishes \emph{on} $Q(1)$ --- this is correct, and is the single most
confusing line in the paper.} \Ledger{D2} $u_0$ is \emph{boundary data}. That
$u_0\equiv0$ on $Q(1)$ does \emph{not} make $\nd u_0$ vanish there: the normal
derivative of the harmonic extension at a boundary point is determined by $u_0$
everywhere, not just nearby, and is in general nonzero on $Q(1)$. It is precisely
this nonzero normal derivative that the recursion drives to $0$. And, per \S1, take
$u_0\not\equiv0$, so that $f\not\equiv0$.
\end{filled}

\begin{ournotation}
\textbf{The order in which everything is chosen} \Ledger{D4} --- the paper gives
only part of this, and the argument is unreadable without it:
\[
 \begin{aligned}
  d,\;p \;&\longrightarrow\; \eta_1\ (\text{Lemma 1})
  \;\longrightarrow\; N \;\longrightarrow\; \tau,\ \beta=\beta(N)\\
  \;&\longrightarrow\; A \;\longrightarrow\; \{K_n\},\{\varepsilon_n\}
  \;\longrightarrow\; \delta_{n+1}\ \text{at each stage.}
 \end{aligned}
\]
A constant written $C$ may depend on $d,p,N$ but never on
$A,K_n,\varepsilon_n,\delta_n$ or $n$. Crucially $\delta_{n+1}$ is chosen
\emph{after} $u_n$ exists, so it may depend on the modulus of continuity of
$\nd u_n$ --- used at D10 and D18 below.
\end{ournotation}

We drop the $\widehat{\phantom{x}}$ and identify functions on $\Rdm$ with their
harmonic extensions.

\subsection*{4.1 The inductive data}

If the $n$th stage has been done $(n\ge0)$ we have: $\delta_n>0$ with
$\delta_n^{-1}\in\mathbb Z$; a subset $\mathcal Q_n\subset\mathcal H_n$, where
$\mathcal H_n$ is the collection of $\delta_n^{-(d-1)}$ cubes of side $\delta_n$
whose union is $Q(1)$; and a smooth $u_n$ with
\begin{equation}\label{eq:IH}
 \Bigl(\int_{V_n}\bigl|\nd u_n\bigr|^p\Bigr)^{1/p}\le A\,e^{-\beta n},
 \qquad V_n=\bigcup\{Q:Q\in\mathcal Q_n\},
\end{equation}
$A$ a large constant independent of $n$. To start, $\delta_0=1$,
$\mathcal Q_0=\{Q(1)\}$.

\begin{expansion}
\textbf{Base case.} \Ledger{D3} At $n=0$, $V_0=Q(1)$ and \eqref{eq:IH} reads
$(\int_{Q(1)}|\nd u_0|^p)^{1/p}\le A$. Since $u_0$ is smooth and $Q(1)$ has measure
$1$, $\nd u_0$ is bounded on $Q(1)$ and the left side is finite; take
$A\ge\|\nd u_0\|_{L^\infty(Q(1))}$. This is one of the constraints on $A$; \S6
imposes the other, and both are lower bounds, so a single large $A$ serves.
\end{expansion}

\subsection*{4.2 Stage $n+1$}

Choose $\delta_{n+1}$ with $\delta_n/\delta_{n+1}\in\mathbb Z$, extremely small
(how small is determined later). Then $\mathcal Q_{n+1}$ is all cubes
$Q\in\mathcal H_{n+1}$ such that
\begin{itemize}[leftmargin=2.4em]
\item[(a)] the cube $Q'\in\mathcal H_n$ with $Q\subset Q'$ satisfies
  $Q'\in\mathcal Q_n$, and
\item[(b)] $\displaystyle\Bigl(\frac1{|Q|}\int_Q\bigl|\nd u_n\bigr|^p\Bigr)^{1/p}
  <K_{n+1}e^{-\beta n}$.
\end{itemize}
Let $a_Q$ be the centre of $Q$ and define
\begin{equation}\label{eq:un1}
 u_{n+1}(x)=u_n(x)+\sum_{Q\in\mathcal Q_{n+1}}\nd u_n(a_Q)\,
   F_{\varepsilon_{n+1}}\bigl(N\delta_{n+1}^{-1}(x-a_Q)\bigr)\cdot
   \frac{\delta_{n+1}}N,
\end{equation}
so that
\begin{equation}\label{eq:gradun1}
 \nabla u_{n+1}(x)=\nabla u_n(x)+\sum_{Q\in\mathcal Q_{n+1}}\nd u_n(a_Q)\,
   \nabla F_{\varepsilon_{n+1}}\bigl(N\delta_{n+1}^{-1}(x-a_Q)\bigr).
\end{equation}

\begin{expansion}
\textbf{Why the factor $\delta_{n+1}/N$ disappears.} \Ledger{D5} Writing
$\Phi_Q(x)=N\delta_{n+1}^{-1}(x-a_Q)$, the chain rule gives
$\nabla\bigl[F(\Phi_Q(x))\bigr]=N\delta_{n+1}^{-1}(\nabla F)(\Phi_Q(x))$, and the
prefactor $\delta_{n+1}/N$ in \eqref{eq:un1} cancels it exactly. This is the reason
for that factor: the \emph{amplitude} of the correction is $O(\delta_{n+1}/N)$, so
$u$ barely moves, while its \emph{gradient} is $O(1)$ --- which is what has to be
$O(1)$ for Lemma 2 to bite. Without the factor the corrections would not converge in
$C^0$; with a larger factor they would not correct the gradient at all.
\end{expansion}

\subsection*{4.3 Support of the correction}

\begin{expansion}
\Ledger{D6, D29} $\supp F_\varepsilon\subset D(0,\varepsilon^{1/2})$, so the
$Q$-term of \eqref{eq:un1} is supported where
$|N\delta_{n+1}^{-1}(x-a_Q)|<\varepsilon_{n+1}^{1/2}$, i.e. in
\[D\Bigl(a_Q,\ \frac{\delta_{n+1}\varepsilon_{n+1}^{1/2}}{N}\Bigr).\]
The original writes the radius as $\delta_{n+1}\varepsilon_{n+1}^{1/2}$, omitting
the $1/N$; since $N$ is fixed this only changes constants, and we keep the correct
radius. For $\delta_{n+1}$ small this disc is contained in $Q$, hence in $Q(1)$.
Consequently
\[u_{n+1}=u_n\ \text{ on }\Rdm\setminus Q(1)\quad\text{for every }n,
  \qquad\text{so}\qquad f=u_0\ \text{ on }\Rdm\setminus Q(1),\]
which is the fact \S1 needs for $f\not\equiv0$.
\end{expansion}

% =====================================================================
\section*{\S5. Lemma 3 and its corollary: controlling the interference}
% =====================================================================

\begin{lemC}
For all $\rho>\mathrm{const}\cdot\delta_{n+1}$ and $x\in\Rdm$,
\begin{equation}\label{eq:L3}
 \sum_{\substack{Q\in\mathcal Q_{n+1}\\|x-a_Q|>\rho}}
  \bigl|\nabla F_{\varepsilon_{n+1}}\bigl(N\delta_{n+1}^{-1}(x-a_Q)\bigr)\bigr|
 \le C\,N^{-d}\,\delta_{n+1}\,\rho^{-1},
\end{equation}
where $C$ depends on a lower bound for $\rho/\delta_{n+1}$.
\end{lemC}

\begin{proof}
The left hand side is
$\le\sum_{|x-a_Q|>\rho}\bigl(N\delta_{n+1}^{-1}|x-a_Q|\bigr)^{-d}$ by the last part
of Lemma 1. Lemma 3 follows by considering
$\sum\delta_{n+1}^{d-1}|x-a_Q|^{-d}$ as a Riemann sum for
$\int_{|x|>\rho}|x|^{-d}$, which is justified as long as $\rho/\delta_{n+1}$ does
not approach $0$.
\end{proof}

\begin{expansion}
\textbf{The arithmetic.} \Ledger{D7} By \eqref{eq:L1grad} applied to
$\varepsilon=\varepsilon_{n+1}$ at the point $y=N\delta_{n+1}^{-1}(x-a_Q)$,
$|\nabla F_{\varepsilon_{n+1}}(y)|\lesssim|y|^{-d}
 =N^{-d}\delta_{n+1}^{d}|x-a_Q|^{-d}$; the other entry $\varepsilon^{-d}$ of the
minimum is not used here because $|x-a_Q|>\rho$ keeps $|y|$ large. Pulling out the
constants,
\[\text{LHS}\ \lesssim\ N^{-d}\delta_{n+1}^{d}\!\!\sum_{|x-a_Q|>\rho}\!\!|x-a_Q|^{-d}
 \;=\;N^{-d}\,\delta_{n+1}\cdot
 \Bigl[\delta_{n+1}^{\,d-1}\!\!\sum_{|x-a_Q|>\rho}\!\!|x-a_Q|^{-d}\Bigr].\]
The bracket is a Riemann sum: the centres $a_Q$ of the cubes of $\mathcal H_{n+1}$
form a lattice of spacing $\delta_{n+1}$ in $\Rdm$, so each contributes a cell of
volume $\delta_{n+1}^{d-1}$, and the bracket approximates
\[\int_{|y|>\rho}|y|^{-d}\,dy
 =\omega_{d-2}\int_\rho^\infty r^{-d}r^{d-2}\,dr
 =\omega_{d-2}\int_\rho^\infty r^{-2}\,dr=\omega_{d-2}\,\rho^{-1}\]
in $\Rdm$. Combining gives \eqref{eq:L3}. Note where $d\ge2$ is needed for the
integral to converge at infinity --- it always does, since the exponent is $-2$
independently of $d$.
\end{expansion}

\begin{filled}
\textbf{The Riemann-sum comparison, made quantitative.} \Ledger{D8} Let
$L=\rho/\delta_{n+1}$, assumed $\ge L_0$ for some fixed $L_0>1$ --- this is the
``$\rho>\mathrm{const}\cdot\delta_{n+1}$'' and the ``$C$ depends on a lower bound
for $\rho/\delta_{n+1}$'' of the statement. For a lattice point $a_Q$ with
$|x-a_Q|>\rho$, every $y$ in its cell $a_Q+[-\frac{\delta_{n+1}}2,
\frac{\delta_{n+1}}2]^{d-1}$ satisfies
$|x-y|\ge|x-a_Q|-\tfrac{\sqrt{d-1}}2\delta_{n+1}
 \ge|x-a_Q|\bigl(1-\tfrac{\sqrt{d-1}}{2L}\bigr)$.
Choosing $L_0=\sqrt{d-1}$ gives $|x-y|\ge\frac12|x-a_Q|$, hence
$|x-a_Q|^{-d}\le2^{d}|x-y|^{-d}$ and
\[\delta_{n+1}^{\,d-1}\!\!\sum_{|x-a_Q|>\rho}\!\!|x-a_Q|^{-d}
 \le2^{d}\!\!\int_{|x-y|>\rho/2}\!\!|x-y|^{-d}dy
 =2^{d}\,\omega_{d-2}\,(\rho/2)^{-1},\]
the cells being disjoint and contained in $\{|x-y|>\rho/2\}$. So
$C=2^{d+1}\omega_{d-2}$ works, and the dependence on the lower bound $L_0$ is
explicit. Without $L\ge L_0$ the cell containing $x$ itself could be included and
the sum would diverge --- which is exactly the caveat the paper states.
\end{filled}

\begin{corC}
If $\delta_{n+1}$ is small enough, then on $\Rdm$,
\begin{equation}\label{eq:cor}
 |\nabla u_{n+1}-\nabla u_n|\le C\,K_{n+1}\,e^{-\beta n}\,\varepsilon_{n+1}^{-d}.
\end{equation}
\end{corC}

\begin{proof}
Assume for notational purposes $x\in Q\in\mathcal Q_{n+1}$ --- the other case is a
little simpler. If $\delta_{n+1}$ is small, then by (b) and smoothness of $u_n$,
\begin{equation}\label{eq:ptwise}
 \bigl|\nd u_n(a_{Q'})\bigr|\le2K_{n+1}e^{-\beta n}
 \qquad\text{for all }Q'\in\mathcal Q_{n+1}.
\end{equation}
So, by \eqref{eq:gradun1},
\begin{align*}
 |\nabla u_{n+1}(x)-\nabla u_n(x)|
 &\le\sum_{\substack{Q'\in\mathcal Q_{n+1}\\Q'\ne Q}}
   \bigl|\nd u_n(a_{Q'})\bigr|\,
   \bigl|\nabla F_{\varepsilon_{n+1}}(N\delta_{n+1}^{-1}(x-a_{Q'}))\bigr|\\
 &\qquad+\bigl|\nd u_n(a_{Q})\bigr|\,
   \bigl|\nabla F_{\varepsilon_{n+1}}(N\delta_{n+1}^{-1}(x-a_{Q}))\bigr|\\
 &\le2K_{n+1}e^{-\beta n}\bigl(C+C\varepsilon_{n+1}^{-d}\bigr),
\end{align*}
the sum by Lemma 3 and the remaining term by Lemma 1.
\end{proof}

\begin{expansion}
\textbf{Two things to fix and one to justify.}
\begin{itemize}[leftmargin=1.4em,itemsep=3pt]
\item \Ledger{D9} The original writes $|\nabla u_n(a_{Q'})|$ throughout this proof.
  But the coefficient in \eqref{eq:gradun1} is $\nd u_n(a_{Q'})$, the \emph{normal}
  derivative, and hypothesis (b) controls only the normal derivative. We have
  written $\nd u_n(a_{Q'})$. (Since $|\nd u_n|\le|\nabla u_n|$, the printed version
  is an over-estimate of the coefficient, but it is not one that (b) supports.)
\item \Ledger{D10} \eqref{eq:ptwise} is the step ``by (b) and smoothness of $u_n$''.
  Hypothesis (b) is an $L^p$ average over $Q'$; to convert it to the pointwise value
  at $a_{Q'}$, note that $\nd u_n$ is continuous, so
  $|\nd u_n(a_{Q'})|^p\le\frac1{|Q'|}\int_{Q'}|\nd u_n|^p+\omega(\delta_{n+1})$
  where $\omega$ is a modulus of continuity for $|\nd u_n|^p$ on $Q(1)$. Choosing
  $\delta_{n+1}$ small enough that
  $\omega(\delta_{n+1})\le(2^p-1)K_{n+1}^pe^{-p\beta n}$ gives
  $|\nd u_n(a_{Q'})|\le2K_{n+1}e^{-\beta n}$. This is legitimate because
  $\delta_{n+1}$ is chosen after $u_n$ (see the ordering block in \S4); $\omega$
  depends on $n$, and that is allowed.
\item \Ledger{D11} The two terms. For $Q'\ne Q$ apply Lemma 3 with
  $\rho=\frac12\delta_{n+1}$: distinct cube centres are $\ge\delta_{n+1}$ apart, so
  every $Q'\ne Q$ has $|x-a_{Q'}|\ge\frac12\delta_{n+1}$ when $x\in Q$; the sum is
  then $\le CN^{-d}\delta_{n+1}(\frac12\delta_{n+1})^{-1}=2CN^{-d}\le C$. For
  $Q'=Q$ we cannot use Lemma 3 ($x$ may equal $a_Q$), and instead use the other
  entry of the minimum in \eqref{eq:L1grad}:
  $|\nabla F_{\varepsilon_{n+1}}|\lesssim\varepsilon_{n+1}^{-d}$. That single term
  is the sole source of the factor $\varepsilon_{n+1}^{-d}$ in \eqref{eq:cor}, and
  hence the sole reason condition (1) of \S7 carries $\varepsilon_{n+1}^{-d}$.
\end{itemize}
\end{expansion}

% =====================================================================
\section*{\S6. Lemma 4: the induction step}
% =====================================================================

\begin{lemD}
If $A$ is large then
$\displaystyle\Bigl(\int_{V_{n+1}}\bigl|\nd u_{n+1}\bigr|^p\Bigr)^{1/p}
 <A\,e^{-\beta(n+1)}$.
\end{lemD}

\subsection*{6.1 The interference estimate $(*)$}

\noindent First we claim: given $\gamma>0$, if $\delta_{n+1}$ is small enough then
for any $x\in Q\in\mathcal Q_{n+1}$
\begin{equation}\label{eq:star2}\tag{$*$}
 \sum_{\substack{Q'\in\mathcal Q_{n+1}\\Q'\ne Q}}
  \bigl|\nd u_n(a_{Q'})\bigr|\,
  \bigl|\nabla F_{\varepsilon_{n+1}}(N\delta_{n+1}^{-1}(x-a_{Q'}))\bigr|
 \le C\,N^{-d}\bigl|\nd u_n(x)\bigr|+\gamma .
\end{equation}

\begin{filled}
\textbf{Order of choices in $(*)$.} \Ledger{D14} The printed proof ``fixes''
$M<\infty$, then at the end says $M$ is arbitrary. Read it as: given $\gamma>0$,
first choose $M$ so large that the far part is $<\gamma/2$; then choose
$\delta_{n+1}$ so small that the near part is within $\gamma/2$ of its limit. Both
choices are made after $u_n$, $K_{n+1}$, $\varepsilon_{n+1}$ are known.
\end{filled}

\begin{proof}[Proof of $(*)$]
Fix $M<\infty$ and split the sum at $|x-a_{Q'}|=M\delta_{n+1}$.

\emph{Near part.}
\begin{align*}
 \sum_{|x-a_{Q'}|<M\delta_{n+1}}
 &\le\sup\Bigl\{\bigl|\nd u_n(a_{Q'})\bigr|:|x-a_{Q'}|<M\delta_{n+1}\Bigr\}
   \cdot\!\!\sum_{Q'\ne Q}\!\!
   \bigl|\nabla F_{\varepsilon_{n+1}}(N\delta_{n+1}^{-1}(x-a_{Q'}))\bigr|\\
 &<CN^{-d}\sup\Bigl\{\bigl|\nd u_n(a_{Q'})\bigr|:|x-a_{Q'}|<M\delta_{n+1}\Bigr\}
 \qquad\text{(Lemma 3 with }\rho=\tfrac12\delta_{n+1})\\
 &<CN^{-d}\Bigl(\bigl|\nd u_n(x)\bigr|+\gamma\Bigr),
\end{align*}
where $\gamma\to0$ as $\delta_{n+1}\to0$ by continuity of $|\nd u_n|$.

\emph{Far part.}
\begin{align*}
 \sum_{|x-a_{Q'}|>M\delta_{n+1}}
 &\le\sup\bigl\{|\nd u_n(y)|:y\in V_{n+1}\bigr\}\cdot\!\!
   \sum_{|x-a_{Q'}|>M\delta_{n+1}}\!\!
   \bigl|\nabla F_{\varepsilon_{n+1}}(N\delta_{n+1}^{-1}(x-a_{Q'}))\bigr|\\
 &\le2K_{n+1}e^{-\beta n}\!\!\sum_{|x-a_{Q'}|>M\delta_{n+1}}\!\!
   \bigl|\nabla F_{\varepsilon_{n+1}}(\cdots)\bigr|
 \qquad(\delta_{n+1}\text{ small enough})\\
 &\le C\,K_{n+1}\,e^{-\beta n}\,M^{-1}
 \qquad\text{(Lemma 3 with }\rho=M\delta_{n+1}).
\end{align*}
Since $M$ is arbitrary we can make this as small as we want, which proves $(*)$.
\end{proof}

\begin{expansion}
\Ledger{D12, D13} Two details. The near part uses Lemma 3 with
$\rho=\frac12\delta_{n+1}$, giving $CN^{-d}\delta_{n+1}\cdot2\delta_{n+1}^{-1}
 =2CN^{-d}$; the sup is then handled by continuity, as in D10. The original writes
``$\gamma\to0$ as $\delta_n\to0$''; it is $\delta_{n+1}$, the side of the
stage-$(n{+}1)$ cubes, that shrinks and forces $a_{Q'}\to x$. The far part uses
Lemma 3 with $\rho=M\delta_{n+1}$, giving
$CN^{-d}\delta_{n+1}(M\delta_{n+1})^{-1}=CN^{-d}M^{-1}$, and the pointwise bound
\eqref{eq:ptwise} for the coefficients.
\end{expansion}

\subsection*{6.2 Type 1 and type 2 cubes}

Call $Q\in\mathcal Q_{n+1}$ \emph{type 1} if
$|\nd u_n(a_Q)|>e^{-4\beta(n+1)}$, and \emph{type 2} otherwise. Fix
$Q\in\mathcal Q_{n+1}$ and write, for $x\in Q$,
\begin{equation}\label{eq:decomp}
 \nd u_{n+1}(x)=\underbrace{\nd u_n(a_Q)}_{\text{the ``}I(a_Q)\text{''}}
 \;+\;\mathrm{Br}(x)\;+\;\nd u_n(a_Q)\,
   \nd F_{\varepsilon_{n+1}}\bigl(N\delta_{n+1}^{-1}(x-a_Q)\bigr),
\end{equation}
where the \emph{bracket} is
\begin{equation}\label{eq:bracket}
 \mathrm{Br}(x)=\nd u_n(x)-\nd u_n(a_Q)
  +\!\!\!\sum_{\substack{Q'\in\mathcal Q_{n+1}\\ Q'\ne Q}}\!\!\!\nd u_n(a_{Q'})\,
   \nd F_{\varepsilon_{n+1}}\bigl(N\delta_{n+1}^{-1}(x-a_{Q'})\bigr).
\end{equation}

\begin{expansion}
\Ledger{D16} The printed bracket is unbalanced (the closing bracket sits after the
sum, and the opening one before $\nd u_n(x)$); \eqref{eq:decomp} is the reading that
makes it an identity, and it is an identity: taking $\nd$ of \eqref{eq:un1} and
adding and subtracting $\nd u_n(a_Q)$ gives exactly these three groups. The
decomposition is engineered so that \emph{$I:=\nd u_n(a_Q)+\text{bracket}$ is nearly
constant on $Q$}, which is what Lemma 2 requires, while the last term is the
correction Lemma 2 acts on.
\end{expansion}

By $(*)$ and smoothness of $\nd u_n$, the bracket may be made less than
$CN^{-d}|\nd u_n(a_Q)|+\gamma$ with $\gamma$ arbitrarily small.

\subsection*{6.3 Type 1 cubes}

For $Q$ type 1 the bracket is therefore $<CN^{-d}|\nd u_n(a_Q)|$ provided $\gamma$
is chosen right. So Lemma 2 applies (using here that $CN^{-d}<cN^{-(d-1)}$ for large
$N$) and we get
\begin{equation}\label{eq:type1}
 \Bigl(\int_Q\bigl|\nd u_{n+1}\bigr|^p\Bigr)^{1/p}
 <e^{-2\beta}\bigl|\nd u_n(a_Q)\bigr||Q|^{1/p}
 <e^{-\frac32\beta}\Bigl(\int_Q\bigl|\nd u_n\bigr|^p\Bigr)^{1/p}
 \quad\text{if }\delta_{n+1}\text{ is small.}
\end{equation}

\begin{expansion}
\textbf{Checking Lemma 2's hypothesis.} \Ledger{D17} Put
$I(x)=\nd u_n(a_Q)+\text{bracket}(x)$, so $I(a_Q)=\nd u_n(a_Q)$ (the bracket
vanishes at $x=a_Q$: its first two terms cancel and the sum is evaluated at
$x=a_Q$ --- strictly, the sum does not vanish at $a_Q$, so $I(a_Q)$ differs from
$\nd u_n(a_Q)$ by at most the same bound, which is absorbed). Then
\[N^{d-1}|I(a_Q)|^{-1}\sup_Q|I-I(a_Q)|
 \le N^{d-1}\cdot\frac{2CN^{-d}|\nd u_n(a_Q)|}{|\nd u_n(a_Q)|}=2CN^{-1},\]
which is $<\tau$ for $N$ large. \emph{This is precisely the parenthetical remark}
``we're using here that $CN^{-d}<cN^{-(d-1)}$ for large $N$'': the gain of one power
of $N$ between the bracket's $N^{-d}$ and the hypothesis's $N^{d-1}$ is what makes
the hypothesis verifiable, and it is why type 1 is defined by a lower bound on
$|\nd u_n(a_Q)|$ --- without one, dividing by $|I(a_Q)|$ is uncontrolled. The $\gamma$
is absorbed because for type 1, $|\nd u_n(a_Q)|>e^{-4\beta(n+1)}$ is a fixed
positive number at this stage, so $\gamma$ may be taken far smaller.

\textbf{The second inequality of \eqref{eq:type1}.} \Ledger{D18} We must pass from
$|\nd u_n(a_Q)||Q|^{1/p}$ to $(\int_Q|\nd u_n|^p)^{1/p}$. By continuity of
$\nd u_n$, for $\delta_{n+1}$ small,
$\frac1{|Q|}\int_Q|\nd u_n|^p\ge|\nd u_n(a_Q)|^p-\omega(\delta_{n+1})
 \ge e^{-\frac12\beta p}|\nd u_n(a_Q)|^p$,
the last step because $|\nd u_n(a_Q)|>e^{-4\beta(n+1)}$ is bounded below at this
stage while $\omega(\delta_{n+1})\to0$. Taking $p$-th roots,
$|\nd u_n(a_Q)||Q|^{1/p}\le e^{\frac12\beta}(\int_Q|\nd u_n|^p)^{1/p}$, and
$e^{-2\beta}\cdot e^{\frac12\beta}=e^{-\frac32\beta}$.
\end{expansion}

\subsection*{6.4 Type 2 cubes}

If $Q$ is type 2 the bracket is $<CN^{-d}|\nd u_n(a_Q)|+\gamma$, which may be taken
$<e^{-4\beta(n+1)}$. So
\begin{equation}\label{eq:type2}
 \int_Q\bigl|\nd u_{n+1}\bigr|^p
 \le\bigl|\nd u_{n}(a_Q)\bigr|^p
   \int_Q\Bigl|1+\nd F_{\varepsilon_{n+1}}\bigl(N\delta_{n+1}^{-1}(x-a_Q)\bigr)
   \Bigr|^p+\int_Q[\ \cdot\ ]^p
 \le2\,e^{-4\beta(n+1)p}\,|Q|,
\end{equation}
since (say) the first integral is $<|Q|$ by Lemma 2.

\begin{expansion}
\Ledger{D19} Unpacking: by \eqref{eq:decomp} and Background \ref{bg:sub},
$|\nd u_{n+1}|^p\le|\nd u_n(a_Q)|^p|1+\nd F_{\varepsilon_{n+1}}(\cdots)|^p
 +|\text{bracket}|^p$. For type 2, $|\nd u_n(a_Q)|^p\le e^{-4\beta(n+1)p}$ and the
bracket is $<e^{-4\beta(n+1)}$, so the second integral is
$\le e^{-4\beta(n+1)p}|Q|$. For the first, we need
$\int_Q|1+\nd F_{\varepsilon_{n+1}}(N\delta_{n+1}^{-1}(x-a_Q))|^p<|Q|$: by the
scaling of \S3.1 this equals $\frac{|Q|}{|Q(N)|}\int_{Q(N)}|1+\nd F_\varepsilon|^p$,
and \S3.4 showed $\int_{Q(N)}(|1+\nd F_\varepsilon|^p-1)<-\eta_1/2<0$, i.e.
$\int_{Q(N)}|1+\nd F_\varepsilon|^p<|Q(N)|$. Hence the first integral is $<|Q|$ and
the two together give $2e^{-4\beta(n+1)p}|Q|$. Note this uses Lemma 2 only in the
degenerate case $I\equiv1$, which is why ``say'' appears in the original.
\end{expansion}

\subsection*{6.5 Summing, and the sign that must change}

We conclude
\begin{align}
 \int_{V_{n+1}}\bigl|\nd u_{n+1}\bigr|^p
 &=\int_{\bigcup\{Q\text{ type }1\}}+\int_{\bigcup\{Q\text{ type }2\}}\notag\\
 &\le e^{-\frac32p\beta}\int_{\bigcup\{Q\text{ type }1\}}\bigl|\nd u_n\bigr|^p
   +C\,e^{-4p\beta(n+1)}\!\!\sum_{Q\text{ type }2}\!\!|Q|\notag\\
 &\le A^pe^{-p\beta(n+1)}\cdot e^{-\frac12p\beta}+C\,e^{-4p\beta(n+1)}\notag\\
 &=A^pe^{-p\beta(n+1)}\Bigl(e^{-\frac12p\beta}+CA^{-p}e^{-3p\beta(n+1)}\Bigr).
 \label{eq:final}
\end{align}
The parenthesis term is $<1$ if $A$ was large enough.

\begin{expansion}
\textbf{The sign, corrected.} \Ledger{D21} The original prints
$e^{+\frac12p\beta}$ in the last two lines, and then asserts the parenthesis is
$<1$. It cannot be: $e^{+\frac12p\beta}>1$ already. The recomputation forces the
minus sign. From \eqref{eq:type1}, raising to the $p$-th power and summing over
type 1 cubes,
\[\int_{\bigcup\{Q\text{ type }1\}}\bigl|\nd u_{n+1}\bigr|^p
 \le e^{-\frac32p\beta}\int_{V_n}\bigl|\nd u_n\bigr|^p
 \le e^{-\frac32p\beta}\cdot A^pe^{-p\beta n}\]
by the induction hypothesis \eqref{eq:IH}. Now factor out $A^pe^{-p\beta(n+1)}$:
\[e^{-\frac32p\beta}A^pe^{-p\beta n}
 =A^pe^{-p\beta(n+1)}\cdot e^{-\frac32p\beta+p\beta}
 =A^pe^{-p\beta(n+1)}\,e^{-\frac12p\beta}.\]
So the exponent is $-\frac12p\beta$, and with it the parenthesis is
$e^{-\frac12p\beta}+CA^{-p}e^{-3p\beta(n+1)}$. The first summand is $<1$ by a fixed
margin, and the second is $\le CA^{-p}$, which is $<1-e^{-\frac12p\beta}$ once
$A$ is large --- this is the second constraint on $A$, the first being the base case
of \S4.1. Both are lower bounds on $A$, so one $A$ serves. \hfill$\square$

\Ledger{D20} The step $\sum_{Q\text{ type }2}|Q|\le1$ is just
$\bigcup\mathcal Q_{n+1}\subset Q(1)$ and $|Q(1)|=1$.
\end{expansion}

% =====================================================================
\section*{\S7. The conditions on $K_n,\varepsilon_n$, and the conclusion}
% =====================================================================

The $K_n$ and $\varepsilon_n$ should satisfy
\begin{align}
 &\textbf{(1)}\quad \sum_n\varepsilon_{n+1}^{-d}K_{n+1}e^{-\beta n}<\infty,
 \label{eq:cond1}\\
 &\textbf{(2)}\quad \sum_n\Bigl(K_{n+1}^{-p}+\varepsilon_{n+1}^{\frac{d-1}2}\Bigr)
   \ \text{is sufficiently small.}\label{eq:cond2}
\end{align}
These are clearly compatible: e.g. $\varepsilon_n=C^{-1}n^{-2}$, $K_n=Cn^{2/p}$ for
a suitably large constant $C$.

\begin{expansion}
\textbf{The exponent in (2).} \Ledger{D23} The original prints
$\varepsilon_{n+1}^{d-1}$. The quantity condition (2) is used to bound, in \S7.3, is
$C\sum\varepsilon_{n+1}^{(d-1)/2}$, and since $\varepsilon<1$ we have
$\varepsilon^{d-1}<\varepsilon^{(d-1)/2}$ --- so smallness of
$\sum\varepsilon^{d-1}$ does \emph{not} give smallness of
$\sum\varepsilon^{(d-1)/2}$. The exponent must be $(d-1)/2$, as written above.

\textbf{Compatibility, checked.} \Ledger{D24} With
$\varepsilon_n=C^{-1}n^{-2}$ and $K_n=Cn^{2/p}$:
\begin{itemize}[leftmargin=1.4em,itemsep=2pt]
\item (1): the terms are
  $C^{d}(n+1)^{2d}\cdot C(n+1)^{2/p}\cdot e^{-\beta n}$, which is
  (polynomial)$\times$(geometric) and so summable for every $C$.
\item (2): $\sum K_{n+1}^{-p}=C^{-p}\sum(n+1)^{-2}=C^{-p}(\pi^2/6-1)$, small for
  $C$ large; and
  $\sum\varepsilon_{n+1}^{(d-1)/2}
   =C^{-\frac{d-1}2}\sum(n+1)^{-(d-1)}$, which \emph{converges iff $d-1\ge2$}, and
  is then small for $C$ large.
\end{itemize}
\textbf{So $d\ge3$ enters here.} \Ledger{A6} For $d=2$ the second series is
$\sum(n+1)^{-1}$, divergent, and this route fails --- as it must, since in two
dimensions the theorem is false. We record this as \emph{a} place where the
hypothesis is used, not \emph{the} place: the imported inequality \eqref{eq:AK} may
also require $d\ge3$, and since we do not have [1] we cannot say. Reporting the
$d=2$ obstruction as though it were the only one would overstate what this paper
shows.
\end{expansion}

\subsection*{7.1 Convergence}

Condition (1) implies, by the Corollary to Lemma 3, that the $u_n$ converge in $C^1$
norm on the closed upper half space. Call the limit function $f$.

\begin{expansion}
\Ledger{D25} The original says ``the corollary to lemma 2''. The Corollary in
question is the one stated after Lemma 3 in \S5, whose proof invokes Lemmas 1 and 3.
\end{expansion}

\begin{filled}
\textbf{From the boundary to the closed half space.} \Ledger{D26, D27} The
Corollary \eqref{eq:cor} bounds $|\nabla u_{n+1}-\nabla u_n|$ \emph{on $\Rdm$}. Two
gaps to close.

\emph{(i) $C^0$ convergence.} From \eqref{eq:un1}, the discs $D(a_Q,\cdot)$ being
disjoint,
\[\|u_{n+1}-u_n\|_\infty
 \le\sup_{Q}\bigl|\nd u_n(a_Q)\bigr|\cdot\|F_{\varepsilon_{n+1}}\|_\infty\cdot
   \frac{\delta_{n+1}}N
 \le 2K_{n+1}e^{-\beta n}\|F_{\varepsilon_{n+1}}\|_\infty\frac{\delta_{n+1}}N,\]
using \eqref{eq:ptwise}. Since $\delta_{n+1}$ is at our disposal \emph{after}
$K_{n+1}$ and $\varepsilon_{n+1}$ are fixed, choose it small enough that this is
$\le2^{-n}$; then $\sum\|u_{n+1}-u_n\|_\infty<\infty$ and $u_n\to f$ uniformly on
$\Rdm$. (The paper never estimates $u_{n+1}-u_n$ itself, only its gradient.)

\emph{(ii) Interior.} $v_n:=u_{n+1}-u_n$ is harmonic in $\Rd$, continuous on
$\overline{\Rd}$, and tends to $0$ at infinity (its boundary data is compactly
supported in $Q(1)$). Each $\partial_jv_n$ is harmonic, bounded, and tends to $0$ at
infinity, so by Background \ref{bg:max}
\[\sup_{\Rd}|\nabla v_n|=\sup_{\Rdm}|\nabla v_n|
 \le CK_{n+1}e^{-\beta n}\varepsilon_{n+1}^{-d}.\]
Condition \eqref{eq:cond1} makes $\sum_n\sup_{\Rd}|\nabla v_n|<\infty$, so
$\nabla u_n$ converges uniformly on $\overline{\Rd}$; with (i), $u_n\to f$ in $C^1$
on $\overline{\Rd}$, $f$ is harmonic in $\Rd$ (a locally uniform limit of harmonic
functions), and $f\in C^1(\overline{\Rd})$. This is the $C^1$-up-to-the-boundary
assertion of the Theorem.
\end{filled}

We have to show $f=\nd f=0$ on a set of positive measure.

\subsection*{7.2 Where $f$ moved}

\begin{equation}\label{eq:fmoved}
 \bigl|\{x\in Q(1):f(x)\ne0\}\bigr|
 \le\sum_n\bigl|\{x\in Q(1):u_{n+1}(x)\ne u_n(x)\}\bigr|,
\end{equation}
and $\{x:u_{n+1}(x)\ne u_n(x)\}$ is a union of at most $\delta_{n+1}^{-(d-1)}$
discs each of radius $\delta_{n+1}\varepsilon_{n+1}^{1/2}/N$ (because
$\supp F_\varepsilon\subset D(0,\varepsilon^{1/2})$). So
\begin{equation}\label{eq:fmeasure}
 \bigl|\{x\in Q(1):f(x)\ne0\}\bigr|
 \le C\sum_n\delta_{n+1}^{-(d-1)}
   \Bigl(\frac{\delta_{n+1}\varepsilon_{n+1}^{1/2}}N\Bigr)^{d-1}
 \le C\sum_n\varepsilon_{n+1}^{\frac{d-1}2}.
\end{equation}

\begin{expansion}
\Ledger{D28} $u_0=0$ on $Q(1)$, so $f\ne0$ at a point of $Q(1)$ forces
$u_{n+1}\ne u_n$ there for some $n$; that is \eqref{eq:fmoved}. There are at most
$|\mathcal H_{n+1}|=\delta_{n+1}^{-(d-1)}$ cubes, hence at most that many discs;
each has measure $c_d(\delta_{n+1}\varepsilon_{n+1}^{1/2}/N)^{d-1}$. Multiplying,
the $\delta_{n+1}$ cancels exactly and $N^{-(d-1)}$ is absorbed into $C$, leaving
$\varepsilon_{n+1}^{(d-1)/2}$ --- which is the exponent condition (2) must control,
and the reason for the correction at D23.
\end{expansion}

\subsection*{7.3 Where the normal derivative failed to vanish}

Also $\nd f=0$ on $\bigcap_nV_n$. Thus
$\bigl|\{x\in Q(1):\nd f(x)\ne0\}\bigr|\le\sum_n|V_n\setminus V_{n+1}|$.

\begin{filled}
\textbf{Why $\nd f=0$ on $\bigcap_nV_n$.} \Ledger{D30} The paper asserts this
without argument. Let $x\in\bigcap_nV_n$. The induction hypothesis \eqref{eq:IH}
gives $\int_{V_m}|\nd u_m|^p\le A^pe^{-p\beta m}$ for every $m$; since
$\bigcap_nV_n\subset V_m$,
\[\int_{\bigcap_nV_n}\bigl|\nd u_m\bigr|^p\le A^pe^{-p\beta m}\xrightarrow[m\to\infty]{}0 .\]
By \S7.1, $\nd u_m\to\nd f$ uniformly on $\Rdm$, so
$\int_{\bigcap_nV_n}|\nd u_m|^p\to\int_{\bigcap_nV_n}|\nd f|^p$ (the domain has
finite measure and the integrands converge uniformly). Hence
$\int_{\bigcap_nV_n}|\nd f|^p=0$ and $\nd f=0$ a.e. on $\bigcap_nV_n$ --- which is
all that is needed, the conclusion being about measure.
\end{filled}

If $Q\notin\mathcal Q_{n+1}$ but $Q\subset Q'\in\mathcal Q_n$, then $Q$ fails (b),
i.e.
\[\frac1{|Q|}\int_Q\bigl|\nd u_n\bigr|^p>K_{n+1}^pe^{-p\beta n},
 \qquad\text{i.e.}\qquad
 |Q|<K_{n+1}^{-p}e^{+\beta pn}\int_Q\bigl|\nd u_n\bigr|^p .\]

\begin{expansion}
\Ledger{D31} The original prints $e^{-\beta pn}$ here. Rearranging
$\frac1{|Q|}\int_Q|\nd u_n|^p>K_{n+1}^pe^{-p\beta n}$ gives
$\int_Q|\nd u_n|^p>|Q|K_{n+1}^pe^{-p\beta n}$ and hence
$|Q|<K_{n+1}^{-p}e^{+p\beta n}\int_Q|\nd u_n|^p$: the exponent flips sign when
$e^{-p\beta n}$ moves across. The next display in the original uses $+\beta pn$,
confirming the slip.
\end{expansion}

So
\begin{equation}\label{eq:VnVn1}
 |V_n\setminus V_{n+1}|\le K_{n+1}^{-p}e^{\beta pn}\int_{V_n}\bigl|\nd u_n\bigr|^p
 \le A^pK_{n+1}^{-p},
\end{equation}
and $\bigl|\{x\in Q(1):\nd f(x)\ne0\}\bigr|\le A^p\sum_nK_{n+1}^{-p}$.

\begin{expansion}
\Ledger{D32} Summing the previous display over the cubes $Q$ making up
$V_n\setminus V_{n+1}$ --- they are disjoint and contained in $V_n$ --- gives the first
inequality; the second is \eqref{eq:IH}, $\int_{V_n}|\nd u_n|^p\le A^pe^{-p\beta n}$,
whose exponential exactly cancels the $e^{+\beta pn}$.
\end{expansion}

\subsection*{7.4 The two costs, added}

\begin{equation}\label{eq:total}
 \bigl|\{x\in Q(1):\nd f(x)\ne0\}\bigr|+\bigl|\{x\in Q(1):f(x)\ne0\}\bigr|
 \le C\sum_n\varepsilon_{n+1}^{\frac{d-1}2}+A^p\sum_nK_{n+1}^{-p}<1,
\end{equation}
if (2) holds. So
$E:=Q(1)\cap\{\nd f=0\}\cap\{f=0\}$ must have positive measure.

\begin{expansion}
\Ledger{D33} $|Q(1)|=1$, and \eqref{eq:total} bounds the measure of the union of the
two bad sets by something $<1$; the complement within $Q(1)$ is $E$, so
$|E|\ge1-(\text{that bound})>0$. Note that condition (2) must be small enough to beat
the constants $C$ and $A^p$, both of which are already fixed at this point --- which
is why (2) says ``sufficiently small'' rather than ``finite''.
\end{expansion}

\subsection*{7.5 From $\nd f=0$ to $\nabla f=0$ --- the step the paper omits}

\begin{filled}
\Ledger{D34} The construction has produced $E\subset Q(1)$ with $|E|>0$,
$f|_E=0$ and $\nd f|_E=0$. The Theorem claims $\nabla f=0$ on a set of positive
measure. The remaining step --- absent from the original --- is that the
\emph{tangential} gradient also vanishes.

Write $\nabla f=(\nabla_Tf,\nd f)$ on the boundary. We know $\nd f=0$ a.e. on $E$.
For the tangential part: $f$ restricted to $\Rdm$ is $C^1$ (\S7.1) and vanishes
identically on $E$, so by Background \ref{bg:density}, $\nabla_Tf(x)=0$ at every
density point $x$ of $E$. Almost every point of $E$ is a density point, so
$\nabla_Tf=0$ a.e. on $E$.

Hence $\nabla f=0$ a.e. on $E$, and discarding a null set we obtain a set of
positive measure on which $f$ and $\nabla f$ both vanish. Together with \S1
($f\not\equiv0$, since $f=u_0\not\equiv0$ off $Q(1)$) and \S7.1 ($f$ harmonic, $C^1$
up to the boundary), this is the Theorem. \hfill$\blacksquare$
\end{filled}

% =====================================================================
\section*{\S8. Remarks (theirs)}
% =====================================================================

\begin{rmk}[1]
To minimize technicalities we did not try to obtain a H\"older estimate on the
gradient although this should be possible along the work of [2].
\end{rmk}

\begin{rmk}[2]
The main question remains open, that is, whether it is possible to have harmonic
functions $C^2$ or $C^\infty$ smooth up to the boundary whose gradients vanish on
sets of positive measure. A variant of the previous construction permits to get a
$C^{2+\varepsilon}$ harmonic function $f:\Rd\to\mathbb R$ such that $f$ and $D^2f$
vanish on a common boundary of positive measure.
\end{rmk}

\begin{ournotation}
The original prints ``$C^2$ of $C^\infty$'' for ``or''. \Ledger{E1} The
$C^{2+\varepsilon}$ variant of Remark 2 is asserted with neither proof nor
reference; we reproduce it as an assertion of the authors and do not verify it.
\Ledger{E2}
\end{ournotation}

% =====================================================================
\section*{\S9. Editorial remarks --- everything later than the paper}
% =====================================================================

Quarantined here, per the blueprint: nothing in this section was used above.

\begin{enumerate}[leftmargin=1.6em,itemsep=4pt]
\item \textbf{The references.} [1] is cited as ``private communication (to
  appear)'' and [2] as ``to appear''. We have not identified the published form of
  either from a source we have read, and do not guess. \Ledger{E4, E5} Identifying
  [1] is the single highest-value follow-up for this digestion: it is the only
  input we could not verify.
\item \textbf{Preprint versus journal version.} This digestion follows
  IHES/M/89/31 (September 1989). The published version is Colloq.\ Math.\ 60/61
  (1990), 253--260, which we could not obtain (the publisher's server was
  unreachable). Numbering and text may differ. \Ledger{F4}
\item \textbf{The five slips.} Listed in \S0. All are typographical or
  sign errors; none affects the theorem. Two further points required a real
  repair rather than a correction: the lower bound for $|1+\nd G_\varepsilon|$
  near $|x|=1$ (\S2.6, \S3.3), and the missing final step \S7.5.
\item \textbf{What remains conditional.} The Alexandrov--Kargaev inequality
  \eqref{eq:AK}. With it, the theorem as stated is proved here in full.
\end{enumerate}

\end{document}
